% BEGIN INCLUDED SGA4_English_current_buildable_source_20260906/source/09/drafts/SGA4_IX_LINES_43_278_English.texfrag
\setcounter{pageoriginal}{0}

\setcounter{section}{-1}
\section{引言}\etiquette[0]{IX.sec0}\pageoriginale

注意, 从本篇开始, 与前面各篇以及拓扑空间的经典上同调理论不同,
为了使主要的命题成立, 必须将讨论限制在\emph{挠}层
(而且常常更精确地限制在与剩余特征互素的挠层).

应当指出, 对特征 $p\neq 0$ 的代数闭域上的簇,
实际上不存在类似于 $k=\CC$ 情形下经典理论的、以整数
(甚至实数) 为系数的“合理”上同调理论 (如 Serre 所指出的).
更精确地说, 不存在一个反变函子 $H^1$,
例如定义在特征 $p>0$ 的代数闭域 $K$ 上光滑投射概形的范畴上,
取值于 $\RR$ (或 $\RR$ 的一个子域) 上有限维向量空间的范畴,
“与乘积可交换”, \ie 满足 $H^1(X\times
Y)\xleftarrow{\sim}H^1(X)\times H^1(Y)$,
且当 $X$ 为亏格 1 的连通曲线时有 $\dim H^1(X)=2$.
事实上, 由此将推出: 若 $X$ 为 $K$ 上的阿贝尔簇,
则从 $\End(X)$ 到 $\End(H^1(X))$ 的映射 $u\mapsto u^1$ 是\emph{加性}的,
因而是 $\End(X)$ 的对偶环的一个\emph{表示}.
但在特征 $p$ 下存在椭圆\pageoriginale 曲线 $X$,
其自同态环是 $\ZZ$ 上一个定四元数代数中的极大序 $A$ ([3], p.198),
而这样的代数显然没有在 $\RR$ 上二维向量空间中的表示,
因为 $A\otimes_{\ZZ}\RR$ 是除代数 (四元数代数).

\section{挠层的形式体系}\etiquette[1]{IX.sec1}

设 $T$ 为拓扑斯, $F$ 为 $T$ 上的阿贝尔层.
在 $F$ 中可以乘以 $n$ ($n\in \ZZ$ 为整数),
这个乘法由 (Ab) 中的类似运算诱导.
用 ${}_nF$ 表示这一乘法的核;
因此, 对 $X\in \Ob T$, ${}_nF(X)$ 为 $F(X)$ 中阶整除 $n$ 的截面所成的群.

用 $\PP$ 表示全体素数的集合.

\begin{enonce*}{定义 1.1}\etiquette[1.1]{IX.def1.1}
设 $p$ 为素数的一个集合. 若典范态射
$$
\varinjlim_n {}_nF\longrightarrow F\text{ 是双射}{\quoi},
$$
则称 $F$ 为 $p$-挠层, \sisi{其中}{或} $p$-层,
其中 $n$ 遍历满足 $\ass n\subset p$ 的整数的集合,
即整除 $n$ 的素数均属于 $p$ 的整数的集合.
若 $p=\PP$ 为全体素数的集合, 就称 $F$ 为挠层.
\end{enonce*}

\begin{enonce*}{命题 1.2}\etiquette[1.2]{IX.prop1.2}
\begin{enumerate}
\renewcommand{\theenumi}{\roman{enumi}}
\renewcommand{\labelenumi}{(\theenumi)}
\item $F$ 是 $p$-挠层, 当且仅当 $F$ 是一个取值于 $p$-挠阿贝尔群的预层的相伴层;
可以取
$$
P=\frac{\lim(\text{\rm 预层})}{\ass n\subset p} {}_nF{\quoi}.
$$

\item 若\pageoriginale $F$ 为 $p$-挠层, 且 $X\in \Ob T$
为拟紧对象\footnote[1]{拓扑斯的一个对象 $X$ 称为拟紧的,
若每个覆盖族 $\{X_i\rightarrow X\}$ 都能由一个有限覆盖族加细.
拓扑斯的态射 $u$ 为拟紧的, 若每个拟紧对象的逆像仍拟紧.},
则 $F(X)$ 为 $p$-挠群.

\item 若 $T$ 局部为有限型 {\rm (VI \hyperref[VI.def1.1]{1.1})},
则 $F$ 为 $p$-挠层, 当且仅当对每个拟紧的 $X\in \Ob T$,
$F(X)$ 为 $p$-挠群. 此时对每个拟紧的 $X$ 及每个 $q$, 有
$$
\varinjlim_{\ass n\subset p} H^q(T/X; {}_nF)\xrightarrow{\sim}
H^q(T/X;F)
$$
.

\item 若 $u:T\rightarrow T'$ 为拓扑斯的态射,
$F'$ 为 $T'$ 上的 $p$-挠层, 则逆像 $u^{\ast}F'$ 为 $T$ 上的 $p$-挠层.

\item 若 $u:T\rightarrow T'$ 为拓扑斯的态射,
$T$ 和 $T'$ 局部为有限型, $u$ 拟紧${}^\ast$,
且 $F$ 为 $T$ 上的 $p$-挠层,
则对每个 $q$, 各 $R^qu_{\ast}F$ 为 $T'$ 上的 $p$-挠层.
\end{enumerate}
\end{enonce*}

\noindent
\emph{证明}:
\begin{enumerate}
\renewcommand{\theenumi}{\roman{enumi}}
\renewcommand{\labelenumi}{(\theenumi)}
\item 显然, 若层 $F$ 是 $p$-挠的, 则 $F$ 是预层 $P$ 的相伴层,
其中截面群 $P(X)$ 为 $F(X)$ 的 $p$-挠元素所成的子群.
反之, 设 $F=\underline{a}P$ (II) 为 $P$ 的相伴层,
其中 $P$ 取值于阿贝尔 $p$-挠群的范畴.
由正合列
$$
0\rightarrow {}_nP \rightarrow P\xrightarrow{n}P
$$
得到正合列
$$
0\rightarrow \underline{a}({}_nP)\rightarrow F\xrightarrow{n}F{\quoi},
$$
从而\pageoriginale $\underline{a}({}_nP)={}_nF$.
因为函子 $\underline{a}$ 与归纳极限可交换, 由
$$
\frac{\lim(\text{\rm 预层})}{\ass n\subset p}{}_nP\xrightarrow{\sim}P
$$
得知还有
$$
\varinjlim_{\ass n\subset p} {}_nF\xrightarrow{\sim}F{\quoi}.
$$

\item 置 $P=\underrightarrow{\lim(\text{\rm 预层})}{}_n F$.
则有 $P\subset F$, 所以 $P$ 是分离预层, 因此
$$
\underline{a}P(X)=\varinjlim (\ker(\prod_i P(X_i)\rightrightarrows)){\quoi},
$$
其中, 如常, 极限取遍各覆盖族 $\{X_i\rightarrow X\}$.
但因为 $X$ 拟紧, 只须取有限覆盖族.
对这样的族, $\prod_iP(X_i)$ 为 $p$-挠的, 所以 $\ker$ 也如此.
因而 $\underline{a}P(X)=F(X)$ 为 $p$-挠的.

\item 要验证 $\varinjlim {}_nF\rightarrow F$ 是双射,
只须对一个生成元族中的每个 $X$, 验证它在 $X\in \Ob T$ 处的值为双射.
因此第一个断言归结为 (ii). 第二个断言是下述结果的推论:

\item 由 (i) 得到, 因为 $u^{\ast}(\underline{a}P')=\underline{a}(u\cdot P')$

\item 由 (iii) 容易得到.
\end{enumerate}

我们还需要一个非阿贝尔的变体.
但定义中有一些技术问题 (我不知道它们是否严重),
因此这里只对概形的平展拓扑给出一个定义.

\begin{enonce*}{定义 1.3}\etiquette[1.3]{IX.def1.3}
设\pageoriginale $p$ 为素数的一个集合.
若 $G$ 的每个有限子集都生成一个阶为 $n$ 的有限子群,
其中 $\ass n\subset p$, 则称群 $G$ 为 \emph{ind-p-群}.
这等价于说, 阶为 $n_{\alpha}$、满足 $\ass n_{\alpha}\subset p$ 的有限子群 $G_{\alpha}$
组成一个滤过集, 且 $G=\varinjlim G_{\alpha}$.
若 $p=\PP$, 则称为 \emph{ind-有限群}, 而不称 ind-p-群.
\end{enonce*}

立即可以验证:
\begin{enonce*}{命题 1.4}\etiquette[1.4]{IX.prop1.4}
\begin{enumerate}
\renewcommand{\theenumi}{\roman{enumi}}
\renewcommand{\labelenumi}{(\theenumi)}
\item ind $p$-群的子群以及商群仍是 ind $p$-群.

\item ind-$p$-群的滤过归纳极限为 ind-$p$-群.

\item ind-p-群的有限投射极限为 ind-p-群.
\end{enumerate}
\end{enonce*}

\begin{enonce*}{定义 1.5}\etiquette[1.5]{IX.def1.5}
设 $X$ 为概形. 若对每个平展 $U\rightarrow X$, 其中 $U$ 拟紧,
$F(U)$ 为 ind $p$-群 (\resp ind-有限群),
则称 $X$ 上的群层 $F$ 为 \emph{ind $p$-群层}
(\resp \emph{ind-有限群层}).
\end{enonce*}

\begin{enonce*}{命题 1.6}\etiquette[1.6]{IX.prop1.6}
设 $F$ 为概形 $X$ 上的群层.
\begin{enumerate}
\renewcommand{\theenumi}{\roman{enumi}}
\renewcommand{\labelenumi}{(\theenumi)}
\item $F$ 为 ind-$p$-群层, 当且仅当对 $X$ 的每个几何点 $\xi$,
纤维 $F_{\xi}$ 为 ind-p-群.

\item 若 $f:X\rightarrow X'$ 为态射,
$F'$ 为 $X'$ 上的 ind-p-群层, 则 $f^{\ast}F'$ 为 ind-p-群层.

\item 若 $f:X\rightarrow X'$ 为拟紧态射,
$F$ 为 $X$ 上的 ind-p-群层, 则 $f_{\ast}F$ 为 ind-p-群层.
\end{enumerate}
\end{enonce*}

\noindent
\emph{证明.}
(i) 假定 $F$ 为 ind-p-群层, $\xi$ 为 $X$ 的一个几何点.
则 $F_{\xi}=\varinjlim F(X')$, 其中 $X'$ 遍历一个在 $X$ 上平展的概形的\pageoriginale 伪滤过系
(VIII \hyperref[VIII.prop3.9]{3.9}).
显然, 仿射的 $X'$ (因而拟紧) 组成一个共尾集,
因此 $F_{\xi}$ 是 ind $p$-群的归纳极限,
由\hyperref[IX.prop1.4]{1.4} (ii), 它是 ind-p-群.
反之, 假定对每个 $\xi$, 纤维 $F_{\xi}$ 是 ind-p-群;
设 $U\rightarrow X$ 平展且 $U$ 拟紧.
设 $S\subset F(U)$ 为有限子集.
对 $U$ 的每个几何点 $\xi$, $S$ 在 $F_{\xi}$ 中的像生成一个有限群.
由于有限群是有限表现的, 容易推出存在一个在 $U$ 上平展且带有点 $\xi$ 的 $U_{\xi}$,
使得 $S$ 在 $F(U_{\xi})$ 中生成一个有限群 (\cf VIII 4).
有限多个这样的 $U_{\xi}$ 构成 $U$ 的一个覆盖, 记为 $\{U_1,\ldots,U_n\}$.
则 $S$ 在 $\prod
F(U_i)$ 中的像生成有限子群,
由于 $F(U)\subset \prod
F(U_i)$, 即得结果.

由 (i) 可见断言 (ii) 平凡, 而 (iii) 是直接的.

% END INCLUDED SGA4_English_current_buildable_source_20260906/source/09/drafts/SGA4_IX_LINES_43_278_English.texfrag


% BEGIN INCLUDED SGA4_English_current_buildable_source_20260906/source/09/drafts/SGA4_IX_LINES_279_650_English.texfrag
\section{可构造层}\etiquette[2]{IX.sec2}

\setcounter{subsection}{-1}
\subsection{}\etiquette[2.0]{IX.subsec2.0}
设 $T$ 为拓扑斯. 回忆, 对每个元素 $S\in \Ob (\Ens)$,
对应一个预层 $P$ 的相伴层 $S_T$, 其中对每个 $X\in \Ob (T)$, 有 $P(X)=S$.
表示 $S_T$ 的 $T$ 的对象为
$$
\coprod_{s\in S}e= ``S\times e \mbox{''}{\quoi},
$$
其中 $e$ 为 $T$ 的终对象.
层 $S_T$ 称为值为 $S$ 的\emph{常值层} (IV).
以显然的方式可定义\emph{常值群层、常值 $A$-模层} ($A$ 为环)
以及\emph{常值阿贝尔层}的概念, 它们的值记为 $M$.
若常值层的态射 $f_T:S_T\rightarrow S'_T$ 来自一个集合态射
$f:S\rightarrow S'$, 则称它为常值的.

若存在终对象的一个覆盖 $\{e_i\rightarrow e\}$,
使得 $F$ 在每个 $e_i$ 上都变为常值的, 则称层 $F$ 为\emph{局部常值}的.
类似地\pageoriginale 定义局部常值\emph{群层}或 $A$-\emph{模}层的概念,
并在 $F$ 和 $G$ 为局部常值层时定义\emph{局部常值态射} $f:F\rightarrow G$ 的概念.
最后, 一个局部常值群层或 $A$-模层称为\emph{有限型}的
(\resp \emph{有限表现}的), 若它的局部值是有限型
(\resp 有限表现) 的群或 $A$-模.

\begin{enonce*}{引理 2.1}\etiquette[2.1]{IX.lem2.1}
\begin{enumerate}
\renewcommand{\theenumi}{\roman{enumi}}
\renewcommand{\labelenumi}{(\theenumi)}
\item 设 $f:F\rightarrow G$ 为 $T$ 上局部常值集合层的态射,
并假定 $F$ 的局部值为有限集. 则 $f$ 局部常值.

\item 设 $f:F\rightarrow G$ 为局部常值 $A$-模层的态射,
其中 $F$ 为有限型. 则 $f$ 局部常值, 且 $f$ 的核及余核局部常值.

\item 设 $0\rightarrow F'\rightarrow F\rightarrow F''\rightarrow 0$
为 $A$-模层 (\resp 群层) 的正合列,
其中 $F'$ 和 $F''$ 局部常值, 且 $F''$ 为有限表现的.
则 $F$ 局部常值.
\end{enumerate}
\end{enonce*}

(i) 和 (ii) 的证明. 毫无困难地归结到 $F$ 和 $G$ 为常值的情形.
考虑 (ii): 设 $F=M_T$, $G=N_T$ ($M$, $N$ 为 $A$-模),
并设 $S\subset M$ 为 $M$ 的一个有限生成集.
则 $f$ 由它在 $S_T$ 上的限制决定,
而若 $f|S_T$ 局部常值, 则 $f$ 局部常值.
因此 (ii) 的第一个断言归结为 (i), 第二个断言则是第一个断言的平凡推论.

所以只须证明 (i).
设 $F=S_T$, $G={S'}_T$ 常值, 并仍用 $f$ 表示对象的态射
$f:S\times e\rightarrow S'\times e$ (记号如上).
设 $f_s:e\rightarrow S'\times e$ $(s\in S)$ 为 $f$ 的各分量:
因为 $S$ 有限, 显然只须处理 $f=f_s$ 的情形,
即 $S$ 为单元素集的情形.
设 $X_{s'}\rightarrow e$ $(s'\in S')$ 为使下图
$$
\xymatrix@C=2cm@R=1.5cm{
X_{s'} \ar[r] \ar[d]& e\ar[d]^{i_{s'}}\\
e   \ar[r]^{f}   & S'\times e
}
$$
笛卡儿的态射族,
其中\pageoriginale $i_{s'}:e\rightarrow S'\times e$ 是到第 $s'$ 个分量的“包含”.
族 $\{i_{s'}\}$ 显然覆盖, 因此 $\{X_{s'}\rightarrow e\}$ 也覆盖.
立即验证 (因为拓扑斯中的余积不交)
%% missing reference
$f$ 在 $X_{s'}$ 上变为常值的, 因而 $f$ 局部常值.

(iii) 的证明. 考虑 $A$-模层的情形.
可以假定 $F'$ 和 $F''$ 常值, 且 $F''$ 的值为有限表现的 $A$-模 $M''$.
若 $M''$ 自由, 因而有一个有限基 $(e_i)_{1\leq i\leq n}$,
则各 $e_i$ 定义 $F''$ 的截面, 从而局部可提升为 $F$ 的截面;
这表明 $F''$ 被 $F'$ 的扩张 $F$ 局部分裂,
证明了 $F$ 局部同构于 $F'\times F''$, 因而局部常值.
一般情形下, 选取一个满同态 $\varphi:L''\rightarrow M''$,
其中 $L''$ 为有限型自由模;
$\varphi$ 的核 $R$ 为有限型, 置
$L=\fprod{L''_T}{F}{M''_T}=\fprod{L''_T}{F}{F''}$,
一方面可见 $L$ 是 $L''_T$ 被 $F'$ 的扩张,
由前所述它局部常值;
另一方面有正合列 $0\rightarrow R_T\rightarrow L\rightarrow
F\rightarrow 0$, 由 (ii), 这蕴含 $F$ 局部常值.

从现在起, 只讨论概形的平展拓扑斯.

\begin{enonce*}{引理 2.2}\etiquette[2.2]{IX.lem2.2}
设 $X$ 为概形, $F$ 为 $X$ 上的局部常值层.
则 $F$ 由一个平展 $Y/X$ 表示.
若进一步 $F$ 的纤维有限 (\resp 且非空),
则 $Y$ 为平展覆盖 (\resp 且满射).
\end{enonce*}

证明. 下降 (当 $F$ 的纤维有限时用 \SGA 1 IX 4.1,
一般情形用 \SGA 3 X 5.4).

\begin{enonce*}{定义 2.3}\etiquette[2.3]{IX.def2.3}
一个集合层 (\resp 群层, \resp $A$-模层) 称为可构造的,
若对每个仿射开集 $U\subset X$, 都存在 $U$ 到有限多个可构造
{\rm (\EGA $0_{{\rm III}}$ 9.1.2)} 的既约局部闭子概形 $U_i$ 的并的分解,
使得 $F$ 在每个 $U_i$ 上\pageoriginale 诱导的层局部常值,
且值有限 (\resp 有限, \resp 有限表现\footnote[1]{当不假定 $A$ 为诺特环时,
为同调代数的目的, 这里最好不只要求 $A$-模 $M$ 有限表现,
而要求它有一个由有限型自由 $A$-模构成的左消解.
于是 (\hyperref[IX.def2.3]{2.3}) 中引入的概念应称为:
$F$ 是 1-可构造的 (更一般地, 对每个整数 $n\geq 0$,
以显然的方式可定义 $n$-可构造性).
本篇暂时沿用 (\hyperref[IX.def2.3]{2.3}) 形式的术语;
它主要在 $A$ 为诺特环时是合理的, 在这一情形下它与刚才指出的概念一致.
关于非诺特情形中的有限性概念, 比较 \SGA 6 I.}).
\end{enonce*}

\setcounter{subsection}{3}
\subsubsection{}\etiquette[2.3.1]{IX.subsubsec2.3.1}
关于局部值的有限性假设, 等价于说对每个几何点 $\xi$,
纤维 $F_{\xi}$ 有限 (\resp 有限, \resp 有限表现).
由\hyperref[IX.lem2.1]{2.1} (i) 立即得知,
群层 $F$ 可构造, 当且仅当 $F$ 的底集合层可构造.

必须注意: 若 $F$ 为阿贝尔层,
则它作为群层 (或等价地, 作为集合层) 可构造,
当且仅当它作为 $\ZZ$-模层可构造且\emph{还}具有有限纤维.
因此, $\ZZ_X$ 作为 $\ZZ$-模可构造, 但作为群层不可构造.

\begin{enonce*}{命题 2.4}\etiquette[2.4]{IX.prop2.4}
\begin{enumerate}
\renewcommand{\theenumi}{\roman{enumi}}
\renewcommand{\labelenumi}{(\theenumi)}
\item 设 $X$ 拟紧且拟分离.
则层 \\
(\resp 群层, \resp $A$-模层) 可构造,
当且仅当存在将 $X$ 分解为可构造局部闭子集之并的分解,
使得 $F$ 在每个 $X_i$ 上变为局部常值且有限
(\resp 有限, \resp 有限表现).

\item 要验证 $F$ 可构造, 只须对 $X$ 的一个给定开覆盖中的各 $U_i$,
验证 $F| U_i$ 可构造.

\item 设 $f:Y\rightarrow X$ 为态射.
若 $F$ 可构造, 则 $f^{\ast}F$ 可构造.

\item 可构造层的纤维非空 (\resp 非零) 的点所成\pageoriginale 的集合,
是一个局部可构造子集 {\rm (\EGA $0_{\rm III}$ 9.1.2)}.

\item 设 $X$ 局部诺特. 则 $X$ 上的集合层
(\resp\ldots) $F$ 可构造, 当且仅当对每个 $x\in X$,
$x$ 的闭包 $\bar{x}$ 都存在一个非空开子集 $U$,
使得 $F$ 在 $U$ 上诱导一个值有限 (\resp\ldots) 的局部常值层.
\end{enumerate}
\end{enonce*}

证明. 显然, 若 $f:Y\rightarrow X$ 为态射,
且存在 $X$ 到可构造局部闭子集 $X_i$ 之并的分解,
使得 $F$ 在每个 $X_i$ 上变为局部常值的,
则对 $Y$ 及 $f^{\ast}F$ 也如此 (\cf \EGA IV 1.8.2).
这证明了 (i) 的蕴含 $\Leftarrow$.
现设 $X$ 拟紧且拟分离, $\{U_1,\ldots,U_n\}$ 为 $X$ 的仿射开覆盖,
使得对每个 $i$, $F|U_i$ 可构造.
要找到 (i) 所述的 $X$ 的分解,
只须对 $U_n$ 和 $Y=X-U_n$ (带诱导的既约结构) 找到这样的分解,
因为由“$X$ 拟分离”的假设, $U_n$ 和 $Y$ 在 $X$ 中可构造.
由可构造性的定义, $U_n$ 有这样的分解,
而 $Y$ 是 $n-1$ 个仿射开集 $V_i=Y\cap U_i$ $(i=1,\ldots,n-1)$ 的并,
且 $F| V_i$ 可构造; 因而由对 $n$ 的归纳得到结果 (i).
同一论证证明 (ii); 事实上, 假设蕴含对每个“充分小”的仿射开集 $U$,
$F| U$ 可构造, 而任意仿射开集 $V$ 都可由有限多个这样的开集覆盖.
断言 (ii) 表明可构造性的概念在 $X$ 上是局部的,
(iii)、(iv) 立即得到.
(v) 的蕴含 $\Rightarrow$ 是直接的, 不依赖诺特假设.
若 $X$ 诺特, 则用通常的“诺特归纳法”证明蕴含 $\Leftarrow$.

\begin{enonce*}{命题 2.5}\etiquette[2.5]{IX.prop2.5}
设 $X$ 拟紧且拟分离, $F$ 为 $X$ 上的可构造群层或 $A$-模层.
则 $F$ 存在一个有限滤过, 其相继商形如 $i_{!}G$,
其中 $i:U\rightarrow X$ 是一个可构造局部闭子集的包含,
$G$ 在 $U$ 上局部常值且可构造.
若 $X$ 诺特, 则存在一个这样的滤过, 其中各 $U$ 不可约.
\end{enonce*}

证明.\pageoriginale 由\hyperref[IX.prop2.4]{2.4} (i),
存在将 $X$ 分解为有限多个可构造局部闭子集 $X_i$ 的并的分解,
使得每个 $F|X_i$ 局部常值且有限 (\resp 有限表现).
写 $X_i=U_i\cap \complement V_i$, 其中 $U_i$ 和 $V_i$ 为可构造开集;
令 $N$ 为由 $\{U_i,V_i\}$ 生成的子拓扑 $T$ 中 $X$ 的开集的数目.
对 $N$ 作归纳: 若 $W$ 是 $T$ 的极小非空元素,
显然 $F|W$ 局部常值且有限表现.
设 $i:W\rightarrow X$ 为包含态射, 令 $Y=X-W$.
有正合列
$$
0\longrightarrow i_{!}i^{\ast}F\longrightarrow F\longrightarrow F|
Y\longrightarrow 0{\quoi},
$$
从而归结为 $Y$ 及层 $F| Y$ 的同一断言, 而此时可应用归纳假设.

\begin{enonce*}{命题 2.6}\etiquette[2.6]{IX.prop2.6}
假定 $A$ 为诺特环.
\begin{itemize}
\item[(i)] 可构造 $A$-模层或集合层的有限投射 (\resp 归纳) 极限可构造.
特别地, 可构造 $A$-模层的态射的核、余核及像均可构造.

\item[(i bis)] 可构造集合层 (\resp 群层) 的有限投射极限可构造,
而若 $f:F\rightarrow G$ 为可构造群层的态射,
则核、像及余核 (若 $\Iim(f)$ 为正规子群) 均可构造.
可构造集合层的有限归纳极限可构造.

\item[(ii)] 设 $0\rightarrow M'\rightarrow M\rightarrow
M''\rightarrow 0$ 为 $A$-模层 (\resp 群层) 的正合列,
其中 $M'$, $M''$ 可构造. 则 $M$ 可构造.

\item[(iii)] 设 $M_1\rightarrow M_2\rightarrow M_3\rightarrow
M_4\rightarrow M_5$ 为 $A$-模层的正合列.
若 $M_i$ 对 $i=1,2,4,5$ 可构造, 则对 $i=3$ 也如此.
\end{itemize}
\end{enonce*}

\emph{证明.} (i) 集合与群的断言\pageoriginale 留给读者.
对 $A$-模层, 只须证明可构造 $A$-模层的态射 $f:F\rightarrow G$ 的核与余核可构造.
断言在 $X$ 上是局部的, 因此可设 $X$ 仿射, 从而拟紧且拟分离.
由\hyperref[IX.prop2.4]{2.4} (i), 归结到 $F$ 与 $G$ 局部常值的情形,
再用\hyperref[IX.lem2.1]{2.1} (ii) 及 $A$ 的诺特假设得到结论.
类似地, 用\hyperref[IX.lem2.1]{2.1} (iii) 证明断言 (ii).
断言 (iii) 立即由 (i) 与 (ii) 得到.

\begin{enonce*}{命题 2.7}\etiquette[2.7]{IX.prop2.7}
设 $X$ 为拟紧且拟分离的概形, $F$ 为集合层 (\resp $A$-模层).
$F$ 可构造的充要条件是,
它同构于一对态射 $H\rightrightarrows G$ 的余核,
其中 $H$ 和 $G$ 是由 $X$ 上平展且有限表现的概形所表示的集合层
(\resp $F$ 同构于同态 $𝒜_{V,X}\rightarrow 𝒜_{U,X}$ 的余核,
其中 $U$ 和 $V$ 是 $X$ 上两个平展且有限表现的概形).
\end{enonce*}

%% A_V [and not 𝒜] in 2.9 \XXX

充分性容易由\hyperref[IX.prop2.6]{2.6} (i bis) 得到.

设 $F$ 为可构造集合层.
利用 $X$ 上平展景中无限余积的存在性, 可见能找到满态射 $G\rightarrow F$,
其中 $G$ 由一个在 $X$ 上平展的概形表示,
而且可进一步设为仿射概形 $G_i$ $(i\in I)$ 的余积,
因而在 $X$ 上分离.
对指标集 $I$ 的每个有限子集 $J$, 令 $G_J$ 为各 $i\in J$ 对应的 $G_i$ 的余积,
令 $F_J$ 为它在 $F$ 中的像.
因为 $G_J$ 及 $F$ 可构造, $F_J$ 也如此 (\hyperref[IX.prop2.6]{2.6} (i bis));
所以, $F$ 与 $F_J$ 在位于 $x$ 上方的 $X$ 的几何点处的纤维相等的那些 $x\in X$
所成的集合 $X_J$ 局部可构造 (\hyperref[IX.prop2.4]{2.4} (iv).
因为各 $X_J$ 组成递增族且并为 $X$, 而 $X$ 拟紧,
某个 $X_J$ 等于 $X$ (\EGA IV 1.9.9).
这表明, 将 $G$ 换成 $G_J$ 后, 可设 $G$ 仿射,
因而在 $X$ 上分离, 并且因 $X$ 拟分离而在 $X$ 上拟紧 (\EGA IV 1.2.4).
于是令 $H=\fprod{G}{G}{F}$.
它是 $\fprod{G}{G}{X}$ 的子层, 后者可表,
所以它本身也可表 (VIII \hyperref[VIII.prop6.1]{6.1}).
又因为 $H$ 可构造 (\hyperref[IX.prop2.6]{2.6} (i bis)),\pageoriginale
由前面的论证再次得知它拟紧, 因而 (它在 $X$ 上分离) 在 $X$ 上有限表现.
这证明了\hyperref[IX.prop2.7]{2.7} 的第一个断言,
第二个断言用同一方法证明 (NB 不必要求 $A$ 诺特).
顺便指出, 这个证明还证明了下列推论:

\begin{enonce*}{推论 2.7.1}\etiquette[2.7.1]{IX.cor2.7.1}
设 $X$ 为概形, $F$ 为 $X$ 上平展的概形.
$X$ 上相应的平展层可构造的充要条件是 $F$ 在 $X$ 上有限表现.
\end{enonce*}

命题\hyperref[IX.prop2.7]{2.7} 在这里主要用于推导可构造层的某些取极限性质:

\begin{enonce*}{推论 2.7.2}\etiquette[2.7.2]{IX.cor2.7.2}
设 $X$ 为拟紧且拟分离的概形.
则 $X$ 上每个集合层 (\resp ind-p-群层, \resp $A$-模层)
都是可构造集合层 (\resp\ldots) 的一个滤过归纳系的归纳极限.
\end{enonce*}

若 $F$ 为集合层, 沿用\hyperref[IX.prop2.7]{2.7} 的证明中的记号,
令 $L$ 为对 $(J,H')$ 的集合, 其中 $J$ 为 $I$ 的有限子集,
$H'$ 为 $H_J=\fprod{G_J}{G_J}{F}$ 的拟紧开子集.
以显然的方式给 $L$ 排序.
注意到 $F$ 为各 $F_J$ 的极限,
而每个 $F_J$ 为各 $\Coker(H'\rightarrow G_J)$ 的极限,
其中 $H'$ 遍历 $H_J$ 的拟紧开集;
所以 $F$ 是由 $\ell=(H',J)\in L$ 对应的
$F_{\ell}=\Coker(H'\rightarrow G_J)$ 所成的归纳系的归纳极限,
这就证明了集合的情形.
$A$-模层的情形用本质相同的方法处理,
先将 $F$ 表为一个同态 $𝒜_{V,X}\rightarrow 𝒜_{U,X}$ 的余核,
其中 $U$, $V$ 为 $X$ 上平展且分离的概形.
ind-$\sheaf{p}$-群的情形需要稍加注意.
沿用前面的记号, 用 $L(G_J)$ 表示由集合层 $G_J$ 生成的自由群层,
用 $N_J$ 表示从 $G_J\rightarrow F$ 得到的同态 $L(G_J)\rightarrow F$ 的核,
用 $L$ 表示对 $(J,N')$ 的集合,
其中 $N'$ 是 $N_J$ 的一个“拟紧”的\emph{集合}子层,
\ie 存在满态射 $P\rightarrow N'$, 其中 $P$ 由一个拟紧 $X$-概形表示.
如下给 $L$ 排序: 若 $J\subset J_1$,
且典范\pageoriginale 同态 $L(G_J)\rightarrow L(G_{J_1})$
将 $N'$ 映入 $N'_1$, 就规定 $(J,N')$ 小于 $(J_1, N'_1)$.
立即可见 $L$ 滤过, 并且如下得到一个以 $L$ 为指标的群层归纳系 $F_{\ell}$:
对 $\ell=(J,N')$, 取 $F_{\ell}=$ 群层 $L(G_J)$
对由 $N'$ 生成的不变群子层的商群层.
立即可见 $F$ 同构于各 $F_{\ell}$ 的归纳极限;
还须证明 $L$ 存在一个在 $L$ 中共尾的子集 $L'$, 使得对 $\ell\in L'$,
$F_{\ell}$ 是一个可构造 $\sheaf{p}$-群层.
为此, 只须证明对每个 $J$, 能找到 $N_J$ 的一个拟紧集合子层 $N'$,
它包含给定的拟紧集合子层 $N'_0$,
并使相应的 $F_{\ell}$ 为可构造 $\sheaf{p}$-群层.
先注意, 为此只须 $F_{\ell}$ 的纤维为有限 $\sheaf{p}$-群:
此时 $F_{\ell}$ 自动可构造, 读者容易验证这一点.
另一方面, 读者也可验证, $X$ 中满足如下条件的点 $x$ 所成的集合 $X_{N'}$
局部可构造: $F_{\ell}$ 在 $x$ 上方一个几何点处的纤维是有限 $\sheaf{p}$-群.
而且, $X$ 的子集 $X_{N'}$ 随 $N'$ 递增,
它们的并为整个 $X$:
最后一点由如下事实得到: 当 $\ell=(J,N')$ 而 $J$ 固定时,
各 $F_{\ell}$ 的归纳极限是 $F$ 的群子层 $\bar{F_J}$,
它由拟紧的集合子层 $F_J$ 生成,
所以 $\bar{F_J}$ 的纤维为有限 $\sheaf{p}$-群.
另一方面, 有限生成元的自由群中的每个有限指数子群都是有限生成元的自由群;
所以, 若将它表为其有限子集的滤过并, 其中某个子集就生成整个子群.
由此容易得知, 各

% END INCLUDED SGA4_English_current_buildable_source_20260906/source/09/drafts/SGA4_IX_LINES_279_650_English.texfrag


% BEGIN INCLUDED SGA4_English_current_buildable_source_20260906/source/09/drafts/SGA4_IX_LINES_651_1136_English.texfrag
$X_{N'}$ 的并确为 $X$.
再次使用 (\EGA IV 1.9.9), 得到某个 $X_{N'}$ 等于 $X$,
这就完成了\hyperref[IX.cor2.7.2]{2.7.2} 的证明.

\setcounter{subsection}{7}
\setcounter{subsubsection}{2}
\paragraph{}\etiquette[2.7.2.1]{IX.par2.7.2.1}
注意, 群概形情形下这个稍显费力的证明,
显示了对于纤维未必有限的群层的可构造性条件, 尚缺乏一个方便的形式体系;
对此我们向读者致歉, 如有必要, 请读者自行填补这一缺憾;
聊以慰藉的是, 在 $X$ 诺特时, 我们将在\hyperref[IX.prop2.9]{2.9} (iii) 中给出更简单的证明.

\begin{enonce*}{推论 2.7.3}\etiquette[2.7.3]{IX.cor2.7.3}
设\pageoriginale $X$ 为拟紧且拟分离的概形,
$F$ 为可构造集合层 (\resp 有限 ind-群层, \resp $A$-模层).
则函子 $\Hom(F,G)$ 与 $G$ 的归纳极限可交换,
其中 $G$ 为可变的集合层 (\resp\ldots).
当 $F$ 为 $A$-模层且 $A$ 诺特时,
关于 $A$-模 $G$ 的函子 $\Ext^{i}(X;F,G)$ 也与归纳极限可交换.
\end{enonce*}

这由\hyperref[IX.prop2.7]{2.7} 及 (VII \hyperref[VII.prop3.3]{3.3})
通过通常的论证得到, 细节留给读者.
类似地, 将\hyperref[IX.prop2.7]{2.7} 与 (VII 5) 的结果结合, 容易得到:

\begin{enonce*}{推论 2.7.4}\etiquette[2.7.4]{IX.cor2.7.4}
设 $I$ 为滤过范畴, $i\mapsto
X_i:I^0\rightarrow (\Sch)$ 为一个函子,
它将箭头映为仿射态射, 将对象映为拟紧且拟分离的概形.
对每个概形 $Y$, 令 $F_{c}(Y)$ 为 $Y$ 上可构造集合层
(\resp ind-$\sheaf{p}$-群层, \resp $A$-模层) 的范畴.
则有范畴等价
$$
F_{c}(X)\xleftarrow{\approx}\varinjlim_i F_{c}(X_i){\quoi},
$$
其中 $X=\varprojlim_i X_i$ {\rm (VII 5)},
右边表示\hyperref[VI.def6.3]{VI 6.3} 中定义的纤维范畴的 $\varinjlim$ 范畴.
特别地, $X$ 上每个可构造集合层 (\resp\ldots)
都同构于某个 $X_i$ 上同类层的逆像.
\end{enonce*}

\begin{enonce*}{命题 2.8}\etiquette[2.8]{IX.prop2.8}
设 $f:X\rightarrow Y$ 为局部有限表现的满态射,
$F$ 为 $Y$ 上的集合层 (\resp\ldots).
则 $F$ 可构造, 当且仅当 $f^{\ast}(F)$ 可构造.
\end{enonce*}

我们将用到下列引理:

\begin{enonce*}{引理 2.8.1}\etiquette[2.8.1]{IX.lem2.8.1}
设 $f:X\rightarrow Y$ 为局部有限表现的满态射,
且 $Y$ 拟紧、拟分离.
则存在 $Y$ 的一个有限分划, 各部分为 $Y$ 上有限表现的子概形 $Y_i$
(特别地,\pageoriginale 每个 $Y_i$ 在 $Y$ 中可构造),
且对每个 $i$, 存在有限满态射 $Y''_i\xrightarrow{g_i}Y'_i\xrightarrow{h_i}Y_i$,
其中 $h_i$ 平展, $g_i$ 局部自由
(\ie 除了有限这一条件外, $g_i$ 还平坦且有限表现) 并且为根式态射,
最后还有一个 $Y$-态射 $Y''_i\rightarrow X$
\footnote[1]{这一断言也见于 \EGA IV 17.16.4.}.
\end{enonce*}

利用 $Y$ 拟紧且拟分离, 容易归结到 $Y$ 拟仿射, 从而分离的情形
(使用 $Y$ 的有限仿射开覆盖 $(U_i)_{1\leq i\leq n}$,
并考虑 $Y_i=U_i-U_i\cap \bigcup_{j< i} U_j$),
再由同一论证归结到 $Y$ 仿射的情形.
概形 $X$ 是仿射开子集 $X_i$ 的并, 且 $X_i$ 在 $Y$ 中的像 $T_i$ 可构造
(\EGA IV 1.8.4); 因 $f$ 满射, 各 $T_i$ 的并为 $Y$,
从而 $Y$ 是有限多个 $T_i$ 的并 (\EGA IV 1.9.9).
因此, 将 $X$ 换成相应的各 $X_i$ 的和概形后, 可设 $X$ 仿射.
标准的取极限程序 (\EGA IV 8.1.2 c)) 使我们还可设 $Y$ 诺特.
通常的诺特归纳法结合取极限程序 (\EGA IV 8.1.2 a)),
将问题归结为 $Y$ 是一个域 $k$ 的谱的情形.
此时 $X\neq ∅$, 所以 $X$ 中存在闭点 $x$,
它对应于 $k$ 的有限扩张 $k'$, 后者又是 $k$ 的一个可分扩张 $k'_s$ 的根式扩张.
于是取 $Y'=\Spec(k'_s)$, $Y''=\Spec(k')$,
这就完成了\hyperref[IX.lem2.8.1]{2.8.1} 的证明.

\hyperref[IX.prop2.8]{2.8} 中的必要性平凡 (\hyperref[IX.prop2.4]{2.4} (iii)),
来证明充分性.
由上面的引理, 归结到 $f$ 形如 $gh$ 的情形,
其中 $g$ 和 $h$ 满足引理中关于 $g_i$, $h_i$ 的条件.
因此可依次假定 $f$ 有限且根式, 此情形由 (VIII \hyperref[VIII.thm1.1]{1.1}) 平凡,
或假定 $f$ 有限平展.
又注意到, 使用 $f^{\ast}(F)$ 可构造这一假设及可构造性的定义,
可以假定 (将 $X$ 换成一个有限表现且结构态射满射的 $X$-预概形
$X'=\coprod_i X_i$ 之后) $f^{\ast}(F)$ 已局部常值.
但当 $f$ 平展且满射时, 这蕴含 $F$ 局部常值.
又因 $f$ 满射, 对 $f^{\ast}(F)$ 的纤维的假设对 $F$ 的纤维仍成立,
所以 $F$ 确实可构造,
这就完成了\hyperref[IX.prop2.8]{2.8} 的证明.

\begin{enonce*}{命题 2.9}\etiquette[2.9]{IX.prop2.9}
设\pageoriginale $X$ 为诺特概形, $A$ 为诺特环,
$p$ 为素数的一个集合.
\begin{enumerate}
\renewcommand{\theenumi}{\roman{enumi}}
\renewcommand{\labelenumi}{(\theenumi)}
\item $X$ 上的集合层 (\resp ind $p$-群层, \resp $A$-模层) 的范畴局部诺特,
即它有一个由诺特对象组成的生成元集合.

\item $X$ 上的集合层 (\resp ind $p$-群层, \resp $A$-模层) $F$ 可构造,
当且仅当它诺特.
若 $F$ 为 $A$-模层, 则 $F$ 可构造,
当且仅当它是有限个形如 $A_{U/X}$ 的层的和的商,
其中 $U\rightarrow X$ 平展且有限型
({\rm \hyperref[IV.subsec2.4]{IV.2.4}} 的记号: $A_{U/X}=A_U$).

%% check the reference IV.2.4 [seems incorrect to me] \XXX

\item 集合层 (\resp\ldots) $F$ 的可构造子层组成一个归纳系,
而 $F$ 是这些子层的归纳极限.
\end{enumerate}
\end{enonce*}

证明. 集合的断言留给读者.
先考虑 $A$-模层的情形.
对 (i), 必须找到一个生成元集合, 其元素是诺特对象.
但形如 $A_{U/X}$ 的层, 其中 $U\rightarrow X$ 平展且有限型,
组成一个生成元集合; 容易看出它们可构造,
所以只须证明下列引理:

\begin{enonce*}{引理 2.10}\etiquette[2.10]{IX.lem2.10}
设 $X$ 诺特. 则每个可构造 $A$-模都是诺特的.
\end{enonce*}

引理的证明.

用诺特归纳法, 假定引理对 $X$ 的每个不同于 $X$ 的子概形 $X'$ 成立.
选取 $X$ 的一个非空开子集 $X_0$,
使得 $F$ 在 $X_0$ 上诱导一个局部常值层 $F_0$.

现设 $\{F_i\}$, $i\in \NN$, 为 $F_0$ 的子层的递增序列.
因为对一个在 $X_0$ 的不可约分支上泛性的几何点 $P$,
$F_0$ 在 $P$ 处的纤维有限型,
所以纤维序列 $(F_i)_P$\pageoriginale 稳定, 可设它为常值的.
现设 $Q$ 为 $X_0$ 的一个几何点, 是 $P$ 的一个\emph{特殊化}
(VIII \hyperref[VIII.def7.2]{7.2}).
因为 $F_0$ 在 $X_0$ 上局部常值,
立即可见特殊化态射 $F_{0_Q}\rightarrow F_{0_P}$ (VIII.7) 是双射,
所以对每个 $i$, $(F_i)_Q\rightarrow (F_i)_P$ 为\emph{单射}.

设 $V/X$ 平展且带有点 $P$,
$s_1,\ldots,s_n$ 为生成 $(F_1)_P$ 的截面 $\in F_1(V)$.
则对每个 $i$ 及 $V$ 上方每个为 $P$ 的特殊化的几何点 $Q$,
各 $s_i$ 都生成 $(F_i)_Q$.
因此层序列 $F_i$ 在 $P$ 的像的一个邻域内常值, 将此邻域记为 $X'\neq ∅$.
令 $Y=X-X'$; 因而只须证明各 $F_i| Y$ 的序列稳定,
这由诺特归纳假设得到.

现在对 $A$-模而言, 断言 (ii) 是直接的:
因为各 $A_{U/X}$ 为诺特生成元, 显然有
$F$ 诺特 $\Leftrightarrow F$ 是有限个层 $A_{U/X}$ 的和的商.
但 $A_{U/X}$ 可构造, 所以由\hyperref[IX.prop2.6]{2.6},
$F$ 诺特 $\Rightarrow F$ 可构造;
反向蕴含已经证明 (\hyperref[IX.lem2.10]{2.10}), 这就确立了 (ii).
最后, 对 $A$-模的断言 (iii) 是局部诺特范畴的一个性质 (\cf [1], Ch. IV).

现考虑 ind-$p$-群的情形.
先注意, 可构造群层 $F$ 必为诺特的.
事实上, 因为纤维有限, 可使用与 $A$-模相同的论证.
现在直接证明断言 (iii), 这将蕴含 (i)---事实上,
它将表明可构造层组成一个生成元集合.
而且, 此时任一诺特层必可构造, 由此得到 (ii).

为证明 (iii), 显然只须证明:
一个 ind $p$-群层 $F$ 的由有限多个截面 $s_i\in F(U_i)$
生成的 (群) 子层 $S$ 可构造, 其中 $U_i\rightarrow X$ 平展且有限型
\footnote[1]{回忆这个层的定义:
令 $L_{U/X}$ 为在群范畴中表示函子 $F\mapsto F(U)$ 的层.
截面 $s_i\in F(U_i)$ 诱导态射 $L_{U_i/X}\rightarrow F$,
$S$ 为和态射 $\amalg_i L_{U_i/X}\rightarrow F$ 的像.
它是“包含”各截面 $s_i$ 的最小子层.}.
(因为\pageoriginale 可构造层诺特, 所以它由有限多个截面生成.)
用诺特归纳法, 可以假定对 $X$ 的每个不同于 $X$ 的闭子概形 $Y$, 这都成立.
由截面生成的子层这一概念与层的逆像可交换;
事实上, 令 $L_{U/X}$ 为在群范畴中表示函子 $F\mapsto F(U)$ 的层.
若 $f:Y\rightarrow X$, 则有 $f^{\ast}(L_{U/X})=L_{\fprod{U}{Y}{X}/Y}$.
因为 $S$ 是各 $L_{U_i/X}$ 的和的像, 而 $f^{\ast}$ 与和可交换,
显然 $f^{\ast}S$ 为 $f^{\ast}F$ 中由各截面 $f^{\ast}(s_i)$ 生成的子层.

因此可假定 $S$ 在每个不同于 $X$ 的闭子集 $Y$ 上诱导可构造层,
从而只须 (\hyperref[IX.prop2.4]{2.4} (i))
证明 $S$ 在 $X$ 的某个适当的非空开子集上诱导可构造层.
将 $X$ 换成充分小的开子集后, 可设每个 $U_i$ 是 $X$ 的一个平展覆盖.
只须证明 $S$ 在某个适当的非空开子集上局部常值且纤维有限.
这一断言对于平展拓扑在 $X$ 上是局部的,
所以可设每个 $U_i$ 完全分裂, 记为 $U_i=\amalg_{n_i}X$
($X$ 的 $n_i$ 个副本的和).
设 $s_{ij}$ $(j=1,\ldots,n_i)$ 为 $F(X)$ 的截面,
使得 $s_i$ $(s_{i1},\ldots,s_{in_i})$.
显然, $S$ 正是 $F$ 中由各 $s_{ij}\in F(X)$ 生成的子层.
因为 $F(X)$ 为 ind $p$-群 (\hyperref[IX.def1.5]{1.5}),
各 $s_{ij}$ 生成有限子群\hyperref[IX.def1.3]{1.3}.
因此可以假定各 $s_i$ 的集合是 $F(X)$ 的一个有限子群;
但此时 $S$ 与由各 $s_i$ 生成的\emph{集合}层相同,
后者由\hyperref[IX.prop2.6]{2.6} (i) 作为集合层可构造,
从而作为群层也可构造, \cqfd

在下面两个命题中, 使用特殊化同态
(VIII \hyperref[VIII.subsec7.7]{7.7}) 给出一个层局部常值或可构造的判别准则.

\begin{enonce*}{命题 2.11}\etiquette[2.11]{IX.prop2.11}
设 $X$ 为概形, $F$ 为可构造集合层
(\resp ind-有限群层, \resp $A$-模层);
在集合的情形下,\pageoriginale 还假定 $F$ 的纤维有限.
设 $x\in X$, $\bar{x}$ 为 $x$ 上方的一个几何点.
$F$ 在 $x$ 的一个邻域内局部常值的充要条件是,
对 $\bar{x}$ 的每个泛化几何点 $\bar{x}'$,
以及每个特殊化箭头 $\bar{x'}\rightarrow \bar{x}$
{\rm (VIII \hyperref[VIII.def7.2]{7.2})},
特殊化同态 {\rm (VIII \hyperref[VIII.subsec7.7]{7.7})}
$F_{\bar{x}}\rightarrow F_{\bar{x'}}$ 是同构.
\end{enonce*}

群层的情形为集合层情形的特例, 因此只讨论后者;
模层的情形以本质相同的方式处理.
条件的必要性平凡, 因此只须证明充分性.
因为纤维 $I=F_{\bar{x}}$ 有限,
将 $X$ 换成一个适当的在 $X$ 上平展的概形 $X'$ 后,
可假定存在从常值层 $I_X$ 到 $F$ 的同态 $u:I_X\rightarrow F$,
它在 $\bar{x}$ 处的纤维上诱导同构.
利用两个层的可构造性及\hyperref[IX.prop2.6]{2.6}, 容易看出:
$X$ 中具有如下性质的点 $z$ 所成的集合 $Z$, 是 $X$ 的一个局部可构造子集:
在位于 $z$ 上方的一个几何点 $\bar{z}$ 处的纤维上诱导的同态为双射.
假设立即蕴含它包含 $x$ 的所有泛化点,
因而 (\EGA IV 1.10.1) 它是 $x$ 的一个邻域,
这就证明 $F$ 在 $x$ 的一个邻域内局部常值 (事实上此处甚至常值),\cqfd

\begin{enonce*}{定义 2.12}\etiquette[2.12]{IX.def2.12}
\footnote[1]{此术语由 S. Lubkin 引入.}:
称概形 $X$ 道路连通, 若对 $X$ 的每对几何点 $P$, $Q$,
存在几何点 $P=P_0,\ldots,P_n$; $Q_1,\ldots,Q_n=Q$
及特殊化 {\rm (VIII \hyperref[VIII.def7.2]{7.2})}
$P_i\rightarrow Q_i$ 和 $P_{i-1}\rightarrow Q_i$ $(i=1,\ldots,n)$.
称 $X$ (对于平展拓扑) 局部道路连通,
若对每个平展 $U\rightarrow X$, 存在 $U$ 的一个平展覆盖
$\{U_i\rightarrow U\}$, 使得对每个 $i$, $U_i$ 道路连通.
\end{enonce*}

立即可以验证, 局部诺特预概形局部道路连通.

\begin{enonce*}{命题 2.13}\etiquette[2.13]{IX.prop2.13}
\begin{enumerate}
\renewcommand{\theenumi}{\roman{enumi}}
\renewcommand{\labelenumi}{(\theenumi)}
\item 设\pageoriginale $X$ 局部道路连通,
$F$ 为 $X$ 上的集合层 (\resp 群层, \resp $A$-模层).
假定 $F$ 的纤维有限 (\resp 有限, \resp 有限表现).
则 $F$ 局部常值, 当且仅当对几何点的每个特殊化 $P\rightarrow Q$,
特殊化态射 {\rm (VIII \hyperref[VIII.subsec7.7]{7.7})} $F_Q\rightarrow F_P$ 是双射.

\item 假定 $X$ 局部诺特, 且 $F$ 的纤维有限 (\resp 有限, \resp 有限表现).
若对每个特殊化 $P\rightarrow Q$,
特殊化态射 $F_Q\rightarrow F_P$ 为单射, 则 $F$ 可构造.

\item 假定 $X$ 局部诺特, 且集合层 $F$ 的纤维有限.
令 $c:X\rightarrow \NN$ 为如下函数:
给每个 $x\in X$ 对应 $F$ 在 $x$ 上方的一个几何点 $\bar{x}$ 处的纤维
$F_{\bar{x}}$ 的元素个数.
则 $F$ 可构造, 当且仅当 $c$ 为可构造函数,
\ie 对每个 $n\in \NN$, $c^{-1}(n)$ 可构造.
\end{enumerate}
\end{enonce*}

\emph{证明.}
\begin{enumerate}
\renewcommand{\theenumi}{\roman{enumi}}
\renewcommand{\labelenumi}{(\theenumi)}
\item 设 $P$ 为 $X$ 的一个几何点.
来找到 $P$ 的一个平展邻域, 即一个带有点 $P$ 的平展 $U\rightarrow X$,
使得 $F$ 在 $U$ 上变为常值的.
取 $P$ 的一个平展邻域 $V'$, 使得 $F(V')\rightarrow F_P$ 满射.
由有限性假设, 这样的 $V'$ 存在.
令 $U\rightarrow V'$ 为 $V'$ 中 $P$ 的一个平展邻域, 使得 $U$ 道路连通.
于是 $F(U)\rightarrow F_P$ 仍满射.
选取子集 $S\subset F(U)$, 它双射地映到 $F_P$.
若 $F$ 为群层 (\resp $A$-模层), 容易看出, 必要时改变 $U$ 后,
可以找到这样的集合, 并且它还是 $F(U)$ 的一个子群 (\resp 子模).
这给出态射 $S_U\rightarrow F| U$, 其中 $S_U$ 为值为 $S$ 的常值层.
有
$$
(S_U)_P\xrightarrow{\sim}F_P{\quoi}.
$$
因为\pageoriginale $F$ (\resp $S_U$) 的纤维上的每个特殊化态射都是双射,
又因 $U$ 道路连通, 得到 $S_U\xrightarrow{\sim}F| U$, 由此即得结果.

\item 断言是局部的, 因此可设 $X$ 诺特.
设 $P$ 为 $X$ 的一个极大点上方的几何点.
存在 $P$ 的一个不可约, 因而道路连通的平展邻域 $V\rightarrow X$,
使得 $F(V)$ 满射地映到 $F_P$.
于是它也满射地映到 $V$ 的每个几何点 $Q$ 处的 $F_Q$,
或等价地, 映到 $V$ 在 $X$ 中的像 $U$ 的每个几何点处的相应纤维;
因为这样的点是 $P$ 的特殊化, 且 $F_Q\rightarrow F_P$ 单射.
因此 $U$ 中的每个特殊化态射都是双射, 所以 $F| U$ 局部常值.
用诺特归纳法, 可设 $F| X-U$ 可构造, 由此即得结果.

\item 显然只须处理函数 $c$ 常值的情形.
设 $V\rightarrow X$ 为 $X$ 的一个极大点上方几何点 $P$ 的平展邻域,
使得 $F(V)$ 满射地映到 $F_P$.
取 $F(V)$ 的一个子集 $S$, 使得 $S\xrightarrow{\sim}F_P$.
于是对常值层 $S_V$ 及几何点的每个特殊化 $P\rightarrow Q$, 有交换图
$$
\xymatrix@C=2cm@R=1.25cm{
(S_V)_P   \ar[r]^{\sim}& F_P\\
(S_V)_Q   \ar[r] \ar[u]^{\text{\rotatebox{90}{$\backsim$}}}  & F_Q. \ar[u]
}
$$
因为 $\lvert F_Q\rvert=\lvert F_P\rvert$,
箭头 $F_Q\rightarrow F_P$ 是双射,
所以由 (i), $F$ 在 $P$ 的一个邻域内局部常值.
由诺特归纳法, 证明完成.
\end{enumerate}

下面的结果 (ii) 将在以后起技术作用;
它使我们能将某些验证归结到常值层的情形.

\begin{enonce*}{命题 2.14}\etiquette[2.14]{IX.prop2.14}
设\pageoriginale $X$ 为概形.
\begin{enumerate}
\renewcommand{\theenumi}{\roman{enumi}}
\renewcommand{\labelenumi}{(\theenumi)}
\item 对每个有限且有限表现的态射 $f:X'\rightarrow X$,
以及 $X'$ 上每个可构造集合层 (\resp ind-有限群层, \resp $A$-模层),
$f_{\ast}(F)$ 可构造.

\item 假定 $X$ 拟紧且拟分离, $F$ 为 $X$ 上的可构造集合层 (\resp\ldots).
则可以找到有限态射的有限族 $(p_i:X'_i\rightarrow X_i)$,
对每个 $i$ 找到 $X'_i$ 上的一个常值可构造集合层 (\resp\ldots) $C_i$,
以及单态射
$$
F\hookrightarrow \prod_i p_{i\ast}(C_i){\quoi}.
$$
\end{enumerate}
\end{enonce*}

先证明 (i). 由定义, 可以假定 $X$ 仿射,
再用\hyperref[IX.cor2.7.4]{2.7.4}, 通过标准取极限程序归结到 $X$ 诺特的情形.
再次取极限, 使用\hyperref[IX.cor2.7.4]{2.7.4}
以及\hyperref[IX.prop2.4]{2.4}. (v), 归结到 $X$ 为域的谱的情形.
此时结论归结到关于 $X$ 上方几何点处 $f_{\ast}(F)$ 的纤维的相应断言,
而后者由 (VIII \hyperref[VIII.prop5.5]{5.5}, \hyperref[VIII.rem5.8]{5.8}) 立即得到.

来证明 (ii). 由\hyperref[IX.prop2.4]{2.4} (i),
存在 $X$ 到子概形 $Y_i$ 的分划,
其包含态射 $u_i:Y_i\rightarrow X$ 为有限表现的,
使得 $F$ 在每个 $Y_i$ 上的限制 $F_i$ 为局部常值层.
于是立即可见, 典范同态 $F\rightarrow
u_{i\ast}(F_i)$
定义一个同态 $F\rightarrow \prod_i
u_{i\ast}(F_i)$, 它是单态射.

在上述构造中, 甚至可以安排每个 $F_i$ 不仅局部常值, 而且\emph{等平凡},
\ie 存在满的\emph{有限}平展态射 $q_i:Y'_i\rightarrow Y_i$,
使得 $q^{\ast}_i(F_i)$ 为 $Y'_i$ 上的常值层, 记为 $C_{iY'_i}$.
例如, 这容易由\hyperref[IX.lem2.8.1]{2.8.1} 的定义及
(VIII \hyperref[VIII.thm1.1]{1.1}) 得到.
另一方面,\pageoriginale 由 (\EGA IV 8.12.6) 形式的“主定理”得知,
存在有限态射 $p_i:Z_i\rightarrow X$
以及 $X$-概形的开浸入 $v_i:Y'_i\rightarrow Z_i$.

因为 $F_i$ 嵌入 $q_{i\ast}(q^{\ast}_i(F_i))=q_{i\ast}(C_{iY'_i})$,
所以 $F$ 嵌入各
$u_{i\ast}(q_{i\ast}(C_{iY'_i}))=p_{i\ast}(v_{i\ast}(C_{iY'_i}))$ 的乘积.
若每个 $Z_i$ 正规, 则下面的引理 (\hyperref[IX.lem2.14.1]{2.14.1})
蕴含 $v_{i\ast}(C_{iY'_i})\simeq C_{iZ_i}$,
所以 $F$ 嵌入各 $p_{i\ast}(C_{iZ_i})$ 的乘积, 证明完成.
一般情形中, 引入 $Z_i$ 的正规化 $\bar{Z}_i$,
$Y'_i$ 在 $\bar{Z}_i$ 中的逆像 $\bar{Y}'_i$,
以及浸入 $\bar{v}_i:\bar{Y}'_i\rightarrow \bar{Z}_i$.
立即可见, $\bar{Y}'_i$ 包含 $\bar{Z}_i$ 的极大点,
所以在 $\bar{Z}_i$ 中稠密.
若 $r_i:\bar{Z}_i\rightarrow Z_i$ 与 $s_i:\bar{Y}'_i\rightarrow Y'_i$ 为投影,
由于后者满射, 得知 $C_{iY'_i}$ 嵌入 $s_{i\ast}(C_{i\bar{Y}'_i})$,
所以 $v_{i\ast}(C_{iY'_i})$ 嵌入
$v_{i\ast}(s_{i\ast}(C_{i\bar{Y}'_i}))=r_{i\ast}(\bar{v}_{i\ast}(C_{i\bar{Y}'_i}))$,
由下面的 (\hyperref[IX.lem2.14.1]{2.14.1}), 后者同构于 $r_{i\ast}(C_{i\bar{Z}_i})$.
因此 $F$ 嵌入各
$p_{i\ast}(r_{i\ast}(C_{\bar{Z}_i}))=t_{i\ast}(C_{i\bar{Z}_i})$ 的乘积,
其中 $t_i=p_ir_i$ 是典范投影 $\bar{Z}_i\rightarrow X$.
若这一态射对每个 $i$ 都有限, 则证明完成.
一般情形中, 只能说 $\bar{Z}_i$ 在 $X$ 上\emph{整},
所以形如 $\sheaf{\Spec}(𝒜_i)$,
其中 $𝒜_i$ 为 $𝒪_X$-代数的拟凝聚层, 在 $𝒪_X$ 上整.
使用 (\EGA IV 1.7.7), 可见 $𝒜_i$ 为其在 $X$ 上有限的子代数
$𝒜_{i,j}$ 的归纳极限,
所以 $\bar{Z}_i$ 是在 $X$ 上有限的预概形
$Z_{ij}=\Spec(A_{ij})$ 的投射极限 (VII 5).
应用 (VII \hyperref[VII.cor5.11]{5.11}), 得到
$t_{i\ast}(C_{i\bar{Z}_i})$ 是各 $t_{ij^{\ast}}(C_{iZ_{ij}})$ 的归纳极限,
其中 $t_{ij}:Z_{ij}\rightarrow X$ 为典范投影.
由 (\hyperref[IX.cor2.7.3]{2.7.3}),
同态 $F\rightarrow t_{i\ast}(C_{i\bar{Z}_i})$
因而经某个适当的 $j=j(i)$ 对应的同态
$F\rightarrow t_{ij\ast}(C_{iZ_{ij}})$ 分解.
现在令 $X'_i=Z_{i,j(i)}$, 便得到\hyperref[IX.prop2.14]{2.14} 所述的嵌入.
还须证明下列引理:

\begin{enonce*}{引理 2.14.1}\etiquette[2.14.1]{IX.lem2.14.1}
设 $X$ 为几何单枝概形 (例如, 正规概形),
$i:U\rightarrow X$ 为拟紧开浸入.
则 $U$ 的闭包 $Z$\pageoriginale 带上诱导的既约结构后, 也几何单枝;
更精确地, 对每个 $z\in Z$, 有
$𝒪_{Z,z}\xleftarrow{\sim}𝒪_{X_{\red},z}$.
另一方面, 假定 $i$ 为支配的, 设 $C$ 为集合,
$C_U$ 为由 $C$ 定义的 $U$ 上的常值层.
则自然同态 $C_X\rightarrow i_{\ast}(C_U)$ 为同构.
若将开浸入 $i$ 换成一个包含 $i:x\rightarrow X$,
其中 $x$ 为 $X$ 的极大点, 也有同样的结论.
\end{enonce*}

使用 (\EGA IV 2.3.11), 第一个断言归结到 $X$ 为局部概形且闭点为 $z$ 的情形;
但此时 $X_{\text{ réd}}$ 为整概形且 $U$ 非空,
所以 $Z=X_{\text{ réd}}$, 断言显然成立
(因此, 此处不必假定 $X$ 几何单枝, 只须假定 $X_{\text{réd}}$ 的局部环为整环).
对第二个断言, 逐个纤维考察,
由 (VIII \hyperref[VIII.rem5.3]{5.3}) \loccit 归结到 $X$ 严格局部的情形,
然后归结为证明 $H^0(X,C_X)\rightarrow H^0(U,C_U)$ 是双射.
此时 $X$ 仍不可约, $U$ 为其非空开子集因而不可约,
上述映射的两端均等同于 $C$, 由此得到所需结论.
(这里使用了 (\EGA IV 18.8.14) 形式的“几何单枝”概念,
即 $X_{\text{ réd}}$ 的各环的严格局部化为整环.)
极大点的包含的情形 (在此陈述以便以后引用) 用完全相同的方式处理.

\begin{enonce*}[remark]{注 2.14.2}\etiquette[2.14.2]{IX.rems2.14.2}
当 $X$ 诺特时, 上述证明 (或其直接推论) 表明,
在\hyperref[IX.prop2.14]{2.14} (ii) 中可假定各 $X'_i$ 为整概形;
若进一步假定 $X$ 普遍日本,
\ie 它的各仿射开子集的环普遍日本
(\EGA ${\rm 0_{IV}}$ 23.1.1, IV 7.7.2),
则还可假定各 $X'_i$ 正规.
\end{enonce*}

% END INCLUDED SGA4_English_current_buildable_source_20260906/source/09/drafts/SGA4_IX_LINES_651_1136_English.texfrag


% BEGIN INCLUDED SGA4_English_current_buildable_source_20260906/source/09/drafts/SGA4_IX_LINES_1137_1339_English.texfrag
\section{Kummer 与 Artin--Schreier 理论}\etiquette[3]{IX.sec3}\pageoriginale

\setcounter{subsection}{-1}
\subsection{}\etiquette[3.0]{IX.subsec3.0}
用 $(\GG_m)_X$ 表示由 $X$ 上的乘法群 $\Spec\Oo_X[t,t^{-1}]$ 表示的层.
因此, 对平展 $U\to X$, $(\GG_m)_X$ 在 $U$ 上的截面为
$\Gamma(U,\Oo_U)$ 的可逆元 $\Gamma(U,\Oo_U^*)$.
设 $n\in\NN$.
$(\GG_m)_X$ 上 $n$ 次幂映射的核为“$n$ 次单位根层”, 记为 $(𝛍_n)_X$.
因此有正合列
\begin{equation*}
0\to(𝛍_n)_X\to(\GG_m)_X\xrightarrow{n}(\GG_m)_X{\quoi}.
\tag{3.1}
\end{equation*}\etiquette[3.1]{IX.eq3.1}
若 $n$ 在 $X$ 上可逆 (即 $n$ 与 $X$ 的每个剩余特征互素),
则 $n$ 次幂映射\emph{作为层的态射}为满射:

\begin{enonce*}{Kummer 理论 3.2}\etiquette[3.2]{IX.Thr3.2}
若 $n\in\NN$ 在 $X$ 上可逆, 则有正合列
$$
0\to(𝛍_n)_X\to(\GG_m)_X\xrightarrow{n}(\GG_m)_X\to0.
$$
\end{enonce*}

事实上, 设 $u\in(\GG_m)_X(U)$, 其中 $U\to X$ 平展.
必须找到一个平展覆盖 $\{U_i\to U\}$,
使得 $u$ 在每个 $U_i$ 上变成 $n$ 次幂.
因为 $n$ 在 $U$ 上可逆, 方程
$$
T^n-u=0
$$
在 $U$ 上“可分”, 即
$$
U'=\Spec\Oo_U[T]/(T^n-u)
$$
在 $U$ 上平展.
因为 $U'\to U$ 满射, 它是一个覆盖, 而 $u$ 在 $U'$ 上变成 $n$ 次幂.

指出\pageoriginale 当 $n$ 在 $X$ 上可逆时,
$(𝛍_n)_X$ 为局部常值层, 实际上局部同构于层 $(\ZZ/n\ZZ)_X$,
在一些重要情形下是常值的.
所以 Kummer 理论给出 $X$ 以某些常值层为值的上同调的信息.
按定义,
$$
H^0(X,(\GG_m)_X)=\Gamma(X,\Oo_X^*){\quoi},
$$
而在维数一有:

\begin{enonce*}{定理 3.3}\etiquette[3.3]{IX.thm3.3}
(“Hilbert 第 90 定理”) 有同构
$$
H^1(X,(\GG_m)_X)\xrightarrow{\sim}\Pic X{\quoi},
$$
其中 $\Pic X$ 为 $X$ 上可逆层的同构类所成的群.
\end{enonce*}

\emph{证明.} 这就是说, 典范态射
$$
H^1(X_{\zar},(\GG_m)_X)\longrightarrow
H^1(X_{\text{\rm ét}},(\GG_m)_X)
$$
为双射. 等价地,
$$
R^1\epsilon_*(\GG_m)_X=0,
$$
其中 $\epsilon_*$ 为显然的景态射
$\epsilon:X_{\et}\to X_{\zar}$ (VIII 4) 的正像.
所以归结到 $X$ 仿射的情形.
由于 $H^1$ 可由 \v{C}ech 上同调计算 (V \hyperref[V.subsec2.5]{2.5}),
且对拟紧 $X$, 只须使用拟紧覆盖,
这立即由层的下降理论得到 (见 FGA 190 或 \SGA VIII 1).

\setcounter{subsection}{3}
\setcounter{subsubsection}{0}
\subsubsection{}\etiquette[3.3.1]{IX.subsubsec3.3.1}
假定 $X$ 诺特且没有嵌入分支, 各 $x_j$ 为 $X$ 的极大点.
令 $R_j$ 为 $X$ 在 $x_j$ 处的局部 Artin 环,
$i_j:\Spec R_j\to X$ 为包含.
有\pageoriginale $X_{\text{\rm ét}}$ (分别地, $X_{\zar}$) 上层的典范单射
$$
(\GG_m)_X\longrightarrow
\prod_j i_{j^{\ast}}(\GG_m)_{\Spec R_j}
$$
.
其余核 $D$ 称为 $X_{\text{\rm ét}}$ (分别地, $X_{\zar}$)
上的 \emph{Cartier 除子层}\footnote[1]{在 \EGA IV 21 中,
使用“除子”一词, 而不是“Cartier 除子”.}.

\begin{enonce*}{推论 3.4}\etiquette[3.4]{IX.cor3.4}
设 $X$ 诺特且没有嵌入分支,
$\epsilon:X_{\text{\rm ét}}\to X_{\zar}$ 为典范景态射,
$D$ 为 $X_{\text{\rm ét}}$ 上的 Cartier 除子层.
则 $\epsilon_*D$ 为 $X_{\zar}$ 上的 Cartier 除子层 $D_{\zar}$,
特别地,
$$
H^0(X_{\text{\rm ét}},D)\simeq H^0(X_{\zar},D_{\zar}){\quoi}.
$$
\end{enonce*}

\emph{证明.} 将 $\epsilon_*$ 应用于正合列
$$
0\to(\GG_m)_X\longrightarrow
\prod_j i_{j^{\ast}}(\GG_m)_{\Spec R_j}
\longrightarrow D\to0
$$
得到一个正合列, 因为 $R^1\epsilon_*(\GG_m)_X=0$
(\hyperref[IX.thm3.3]{3.3}); 第一个断言由此得到.

当 $X$ 的特征为 $p\neq0$ 时,
在研究 $p$-挠系数时, 下列结果有时替代\hyperref[IX.Thr3.2]{3.2}:

\begin{enonce*}{Artin--Schreier 理论 3.5}\etiquette[3.5]{IX.Thr3.5}
若 $X$ 的特征为 $p>0$, 则有正合列
$$
0\to(\ZZ/p\ZZ)_X\to\Oo_X\xrightarrow{\mathfrak p}\Oo_X\to0{\quoi},
$$
其中 $\mathfrak p$ 为由 $\mathfrak p(f)=f^p-f$ 定义的加法群态射.
\end{enonce*}

$\mathfrak p$ 的核是 $\Oo_X$ 中局部属于素域的截面所成的层,
因而是常值层.
态射 $\mathfrak p$ 对平展拓扑是\pageoriginale 满射,
因为方程 $T^p-T-a$ ($a\in\Gamma(U,\Oo_U)$) 在特征 $p$ 下总是“可分”的,
即定义 $U$ 的一个平展覆盖.
应用这一正合列时, 当然将使用 VII \hyperref[VII.subsec4.2]{4.2}.

\begin{enonce*}{命题 3.6}\etiquette[3.6]{IX.prop3.6}
\begin{enumerate}
\renewcommand{\theenumi}{\roman{enumi}}
\renewcommand{\labelenumi}{(\theenumi)}
\item 设 $X$ 为概形, $i:x\to X$ 为一个点的包含.
则
$$
H^1(X,i_*(\ZZ)_x)=0,
$$
其中 $(\ZZ)_x$ 为值为 $\ZZ$ 的“常值层”
(见\hyperref[IX.subsec2.0]{2.0}).

\item 若 $X$ 不可约, 且对 $X$ 的每个几何点 $\bar x$,
$X$ 在 $\bar x$ 处的严格局部化 $X_{\bar x}$ 不可约
(应称“$X$ 几何单枝”, 见 \EGA IV 18.8.14), 则
$$
H^1(X,(\ZZ)_X)=0{\quoi}.
$$
\end{enumerate}
\end{enonce*}

\emph{证明.}
\begin{enumerate}
\renewcommand{\theenumi}{\roman{enumi}}
\renewcommand{\labelenumi}{(\theenumi)}
\item Leray 谱序列
$$
E^{pq}_2=H^p(X,R^qi_*(\ZZ)_x)\Longrightarrow
H^{p+q}(x,(\ZZ)_x)
$$
给出单射 $H^1(X,i_*(\ZZ)_x)\to H^1(x,(\ZZ)_x)$.
因为 $x=\Spec k(x)$ 为一个域的谱,
众所周知 $H^1(x,(\ZZ)_x)=0$.
(这由 $H^q(x,F)$ 在 $q>0$ 时总是挠的这一事实
(见 VII \hyperref[VII.rem2.1]{2.2} 及 CG),
以及正合列 $0\to\ZZ\xrightarrow{n}\ZZ\to\ZZ/n\to0$ 得到.)
这证明了 (i).

\item 设 $i:x\to X$ 为 $X$ 的泛点的包含.
在 (ii) 的假设下, 立即可见典范态射 $(\ZZ)_X\to i_*(\ZZ)_x$ 为双射
(反之也成立); 因此 (ii) 归结到 (i).
\end{enumerate}

\begin{enonce*}[remark]{注 3.7}\etiquette[3.7]{IX.rems3.7}
在\hyperref[IX.prop3.6]{3.6} (i) 与 (ii) 中,
$\ZZ$ 可换成任意无挠阿贝尔群.
另一方面, 若去掉 $X$ 几何单枝的假设, (ii) 就不成立;
例如取 $X$ 为域 $K$ 上有一个普通二重点的代数曲线, 即可看出这一点
(计算此时的 $H^1(X,\ZZ)$!).
\end{enonce*}


% END INCLUDED SGA4_English_current_buildable_source_20260906/source/09/drafts/SGA4_IX_LINES_1137_1339_English.texfrag


% BEGIN INCLUDED SGA4_English_current_buildable_source_20260906/source/09/drafts/SGA4_IX_LINES_1340_1459_English.texfrag
\section{代数曲线的情形}\etiquette[4]{IX.sec4}

\setcounter{subsection}{-1}
\subsection{}\etiquette[4.0]{IX.subsec4.0}
设 $X$ 为维数 $1$ 的诺特概形.
令 $R(X)$ 为 $X$ 上的有理函数环, $i:\Spec R(X)\to X$ 为包含.
环 $R(X)$ 为 Artin 环, 是 $X$ 在各极大点处的局部环的乘积,
可以应用 VIII 2 的结果.

设 $F$ 为 $\Spec R(X)$ 上的阿贝尔层,
考虑 $X$ 上的层 $R^qi_*F$, $q>0$.
由 VIII \hyperref[VIII.thm5.2]{5.2} 立即得知,
$R^qi_*F$ 在位于 $X$ 的极大点上方的几何点处的纤维为零.
因为 $X$ 维数为 $1$, 这样的层具有很特殊的形式:

\begin{enonce*}{引理 4.1}\etiquette[4.1]{IX.lem4.1}
设 $X$ 诺特, $F$ 为 $X$ 上的层.
下列条件等价 (满足它们的层称为“摩天层”):
\begin{enumerate}
\renewcommand{\theenumi}{\roman{enumi}}
\renewcommand{\labelenumi}{(\theenumi)}
\item $F$ 在位于非闭点 $x\in X$ 上方的 $X$ 的每个几何点 $\bar x$ 处的纤维为零.

\item 每个截面 $s\in F(X')$, 其中 $X'\to X$ 平展且有限型,
在 $X$ 的有限多个闭点之外为零.

\item 有同构 $F\simeq\bigoplus_x i_{x*}i_x^*F$,
其中 $x$ 遍历 $X$ 的闭点, $i_x:x\to X$ 为包含.
\end{enumerate}
\end{enonce*}

\emph{证明.}\pageoriginale 假定 (i) 成立,
设 $z\in F(X')$, 其中 $X'\to X$ 平展且有限型.
则 $z$ 的支集 (VIII \hyperref[VIII.subsec6.6]{6.6})
是由\emph{闭}点组成的 $X'$ 的闭子集,
由于 $X'$ 诺特, 它有限.
因此 (i) 蕴含 (ii).
(ii) 到 (i) 的蕴含由纤维的计算 (VIII \hyperref[VIII.prop3.9]{3.9}) 得到,
而 (iii) 蕴含 (i), 因为对 $X$ 的闭点 $x$ 上的任意层族 $(G_x)$,
$X$ 上的层 $F=\bigoplus i_{x*}(G_x)$ 满足 (i):
纤维函子与直和可交换, 而 $i_{x*}(G_x)$ 的支集显然包含于 $\{x\}$.

还须证明 (ii) 蕴含 (iii). 考虑典范态射
$$
F\longrightarrow\prod_x i_{x*}i_x^*F{\quoi},
$$
其中 $x$ 遍历 $X$ 的闭点.
由 (ii), 它经
$$
F\longrightarrow\bigoplus_x i_{x*}i_x^*F{\quoi},
$$
分解, 这给出所需的态射. 其双射性的验证留给读者.

\setcounter{subsection}{1}
\subsubsection{}\etiquette[4.1.1]{IX.subsubsec4.1.1}
将引理应用于上面引入的层 $R^qi_*F$. 有
$$
R^qi_*F\xrightarrow{\sim}
\bigoplus_x i_{x*}i_x^*(R^qi_*F),
$$
其中 $x$ 遍历 $X$ 的闭点.
由 VIII 5, $R^qi_*F$ 在 $X$ 的一个几何点 $\bar x$ 处的纤维为
$$
(R^qi_*F)_{\bar x}\simeq H^q(\Spec R(X_{\bar x}),F),
$$
其中 $R(X_{\bar x})\simeq R(X)\otimes_{\Oo_X}\Oo_{X_{\bar x}}$,
$X_{\bar x}$ 为 $X$ 在 $\bar x$ 处的严格局部化 (VIII 4),
$F$ 也表示在 $\Spec R(X_{\bar x})$ 上诱导的层.
现假定 $X$ 的每个闭点 $x$ 的剩余域可分闭.
于是 (另见 VI 1)
\begin{align*}
H^p(X,R^qi_*F)
&\simeq H^p\left(X,\bigoplus_x i_{x*}i_x^*(R^qi_*F)\right)\\
&\simeq \bigoplus_x H^p\left(X,i_{x*}i_x^*(R^qi_*F)\right)\\
&\simeq \bigoplus_x H^p\left(x,i_x^*(R^qi_*F)\right)\\
&\simeq
\begin{cases}
\displaystyle\bigoplus_x H^q(\Spec R(X_x),F)
  &\mbox{若 }p=0,\\
0 &\mbox{若 }p>0{\quoi}.
\end{cases}
\end{align*}
因此:

\begin{enonce*}{推论 4.2}\etiquette[4.2]{IX.cor4.2}
设 $X$ 为维数 $1$ 的诺特概形,
使得对 $X$ 的每个闭点 $x$, 剩余域 $k(x)$ 可分闭.
设 $F$ 为 $\Spec R(X)$ 上的阿贝尔层,
其中 $R(X)$ 为 $X$ 上的有理函数环. 考虑 Leray 谱序列
$$
E^{pq}_2=H^p(X,R^qi_*F)\Longrightarrow
H^{p+q}(\Spec R(X),F){\quoi}.
$$
有
\begin{align*}
E^{pq}_2&=0\quad\mbox{若 }p>0\mbox{ 且 }q>0{\quoi},\\
E^{0q}_2&=\bigoplus_x H^q(\Spec R(X_x),F)
  \quad\mbox{若 }q>0{\quoi},
\end{align*}
其中 $x$ 遍历 $X$ 的闭点, $X_x$ 为 $X$ 在 $x$ 处的严格局部化.

因而\pageoriginale 谱序列给出正合列
\begin{align*}
0\longrightarrow H^1(X,i_*F)&\longrightarrow H^1(\Spec R(X),F)\\
&\longrightarrow\bigoplus_xH^1(\Spec R(X_x),F)\\
&\longrightarrow H^2(X,i_*F)\longrightarrow H^2(\Spec R(X),F)\\
&\longrightarrow\bigoplus_xH^2(\Spec R(X_x),F)\longrightarrow\cdots .
\end{align*}
\end{enonce*}


% END INCLUDED SGA4_English_current_buildable_source_20260906/source/09/drafts/SGA4_IX_LINES_1340_1459_English.texfrag


% BEGIN INCLUDED SGA4_English_current_buildable_source_20260906/source/09/drafts/SGA4_IX_LINES_1460_1606_English.texfrag
\setcounter{subsection}{2}
\subsection{}\etiquette[4.3]{IX.subsec4.3}
再假定对给定素数 $\ell$, $\Spec R(X)$ 的 $\ell$-上同调维数至多为 $1$
(即对 $\Spec R(X)$ 上每个 $\ell$-挠层 $F$,
当 $q>1$ 时有 $H^q(\Spec R(X),F)=0$).
设 $K^1,\ldots,K^m$ 为 Artin 环 $R(X)$ 的各剩余域,
即 $X$ 的各极大点的剩余域.
由 VIII 2 中的对应, $\Spec R(X)$ 的 $\ell$-上同调维数为
$$
\cd_\ell(\Spec R(X))
=\mathop{\sup}_i\{\cd_\ell(G(\overline K^i/K^i))\}{\quoi},
$$
其中 $\cd_\ell(G)$ 表示 pro-有限群 $G$ 的 $\ell$-上同调维数 (见 CG).

对 $X$ 的一个严格局部化 $\bar X(x)$,
$R(\bar X(x))$ 的每个剩余域 $K_x$ 显然是某个 $K^i$ 的可分
(可能无限) 扩张, 因而由 CG II 4.1, 命题 10,
$$
\cd_\ell(G(\bar K_x/K_x))
\leq\cd_\ell(G(\bar K^i/K^i))
$$
.
所以也有 $\cd_\ell(\Spec R(\bar X(x)))\leq1$.

将\hyperref[IX.cor4.2]{4.2} 的正合列应用于 $\ell$-挠层 $F$, 得到
\begin{equation*}
\begin{cases}
0\to H^1(X,i_*F)\to H^1(\Spec R(X),F)\to\\
\qquad\displaystyle\bigoplus_xH^1(\Spec R(\bar X(x)),F)
\to H^2(X,i_*F)\to0{\quoi},\\
H^q(X,i_*F)=0\quad\mbox{当 }q>2{\quoi}.
\end{cases}
\tag{4.3.1}
\end{equation*}\etiquette[4.3.1]{IX.eq4.3.1}

注意\pageoriginale 若 $X$ 为可分闭域上的代数曲线,
则假设 $\cd_\ell(\Spec R(X))\leq1$ 成立
(Tsen 定理 [6, 定理\hyperref[IX.subsec4.3]{4.3}]).
在这种情形下, $\Ogg[2]$ 及 \v{S}afarevi\v{c} 研究过类似的正合列.

\subsection{}\etiquette[4.4]{IX.subsec4.4}
假定 $\ell$ 在 $X$ 上可逆, 取 $F=(\GG_m)_{R(X)}$.
仍在假设 $\cd_\ell(\Spec R(X))\leq1$ 下,
应用\hyperref[IX.thm3.3]{3.3} 及 [CG, 第 I 章, 命题 12], 得到
\begin{equation*}
\begin{cases}
H^1(\Spec R(X),\GG_m)=0,\\
H^q(\Spec R(X),\GG_m)\;\displaystyle{\mathop{=}_{\ell}}\;0
\quad(q\geq2){\quoi},
\end{cases}
\tag{4.4.1}
\end{equation*}\etiquette[4.4.1]{IX.eq4.4.1}
其中等号下的 $\ell$ 表示左边的群
(它总是挠群) 没有 $\ell$-挠元.

将 $R(X)$ 换成 $R(\bar X(x))$, 同样的结果成立.
由于 $\Spec R(X)$ 上的层 $\GG_m$ 显然在 $\Spec R(\bar X(x))$ 上诱导层 $\GG_m$,
\hyperref[IX.cor4.2]{4.2} 给出
\begin{equation*}
\begin{cases}
H^1(X,i_*(\GG_m)_{R(X)})=0,\\
H^q(X,i_*(\GG_m)_{R(X)})
\;\displaystyle{\mathop{=}_{\ell}}\;0
\quad\mbox{若 }q>1{\quoi}.
\end{cases}
\tag{4.4.2}
\end{equation*}\etiquette[4.4.2]{IX.eq4.4.2}
对 $X$ 上的每个层 $F$, \hyperref[IX.lem4.1]{4.1} 表明,
典范同态 $F\to i_*i^*F$ 的核 $P$ 及余核 $Q$ 为摩天层.
由于假定 $X$ 的闭点的剩余域可分闭,
当 $j>0$ 时便有 $H^j(X,P)=H^j(X,Q)=0$.
上同调长正合列立即蕴含
$$
H^j(X,F)\longrightarrow H^j(X,i_*i^*F)
$$
在\pageoriginale $j\geq2$ 时为同构.
考虑\eqhyperref[IX.eq4.4.2]{4.4.2}, 得到
\begin{equation*}
H^j(X,\GG_m)\displaystyle{\mathop{\;=\;}_{\ell}}0
\quad\mbox{当 }j\geq2{\quoi},
\tag{4.5}
\end{equation*}\etiquette[4.5]{IX.eq4.5}
它在 $X$ 是维数 $1$ 的诺特概形,
满足 $\cd_\ell(R(X))\leq1$, 且闭点的剩余域可分闭时成立.
现在由 Kummer 理论 (\hyperref[IX.Thr3.2]{3.2}) 及\hyperref[IX.thm3.3]{3.3} 得到:

\begin{enonce*}{定理 4.6}\etiquette[4.6]{IX.thm4.6}
设 $X$ 诺特且维数为 $1$, $n$ 为在 $X$ 上可逆的整数.
假定对整除 $n$ 的每个素数 $\ell$, 有 $\cd_\ell(R(X))\leq1$,
并且对 $X$ 的每个闭点 $x$, 剩余域 $k(x)$ 可分闭.
则
$$
H^q(X,(𝛍_n)_X)=0\quad\mbox{若 }q>2{\quoi},
$$
而 $q\leq2$ 的上同调由下列正合列给出
\begin{align*}
0\to H^0(X,(𝛍_n)_X)&\to\Gamma(X,\Oo_X^*)
\xrightarrow{n}\Gamma(X,\Oo_X^*)\to\\
H^1(X,(𝛍_n)_X)&\to\Pic X\xrightarrow{n}\Pic X
\to H^2(X,(𝛍_n)_X)\to0{\quoi}.
\end{align*}
\end{enonce*}

\begin{enonce*}{推论 4.7}\etiquette[4.7]{IX.cor4.7}
设 $X$ 为可分闭域 $k$ 上的完备连通曲线, 且其特征与 $n$ 互素.
则
\begin{align*}
H^0(X,(𝛍_n)_X)&\simeq𝛍_n(k){\quoi},\\
H^1(X,(𝛍_n)_X)&\simeq{}_n(\Pic X)
  =\{\text{$\Pic X$ 中阶整除 $n$ 的元素}\},\\
H^2(X,(𝛍_n)_X)&\simeq(\Pic X)/n
  \simeq(\ZZ/n\ZZ)^c{\quoi},
\end{align*}
其中\pageoriginale $c$ 为 $X$ 的不可约分支数.
\end{enonce*}

对最后一个断言, 回忆这样的 $X/\Spec k$ 的 Picard 概形 [TDTE V] 有分解
\begin{equation*}
0\longrightarrow A\longrightarrow\sheaf{\Pic}_{X/k}
\longrightarrow\ZZ^c\longrightarrow0
\tag{4.8}
\end{equation*}\etiquette[4.8]{IX.eq4.8}
其中 $A$ 为 $\Spec k$ 上的连通代数群.
这样的 $A$ 上“乘以 $n$”的映射平展, 因而在 $k$-值点上满射.
于是
$$
{}_n(\Pic X)\simeq{}_n(A(k))
$$
且
$$
(\Pic X)/n\simeq(\ZZ/n\ZZ)^c{\quoi}.
$$

若 $X$ 在 $\Spec k$ 上光滑, 从而其 Jacobian $A$ 是阿贝尔簇,
则群 \({}_n(\Pic X)\) 同构于 $(\ZZ/n)^{2g}$,
其中 $g$ 是 $X$ 的亏格. 当然, 也有 $𝛍_n(k)\simeq\ZZ/n\ZZ$.
因此各上同调群作为 $\ZZ/n$-模的秩为 $1$, $2g$ 及 $1$, 正如所预期的.


% END INCLUDED SGA4_English_current_buildable_source_20260906/source/09/drafts/SGA4_IX_LINES_1460_1606_English.texfrag


% BEGIN INCLUDED SGA4_English_current_buildable_source_20260906/source/09/drafts/SGA4_IX_LINES_1607_1883_English.texfrag
\section{迹方法}\etiquette[5]{IX.sec5}

\subsection{}\etiquette[5.1]{IX.subsec5.1}
设 $X$ 为概形, $f:X'\rightarrow X$ 为平展覆盖, $F$ 为 $X$ 上的层.
于是伴随公式给出\emph{限制}态射
$$
F\xrightarrow{\res}f_{\ast}f^{\ast}F{\quoi}.
$$
因为 $f$ 有限, 有 $H^q(X,f_{\ast}f^{\ast}F)\simeq H^q(X',
f^{\ast}F)$, 从而有一个也称为限制的态射:
\begin{equation*}
H^q(X,F)\xrightarrow{\res}H^q(X',f^{\ast} F){\quoi},\tag{5.1.1}
\end{equation*}\etiquette[5.1.1]{IX.eq5.1.1}
它\pageoriginale 就是显然的态射.
但当 $f$ 平展时, 还有另一个态射, 称为\emph{迹态射},
\begin{equation*}
f_{\ast}f^{\ast} F\xrightarrow{\tr} F{\quoi},\tag{5.1.2}
\end{equation*}\etiquette[5.1.2]{IX.eq5.1.2}
它在上同调上诱导态射
\begin{equation*}
H^q(X', f^{\ast} F)\xrightarrow{\tr}H^q
(X,F){\quoi}.\tag{5.1.3}
\end{equation*}\etiquette[5.1.3]{IX.eq5.1.3}
迹态射由下列两个性质刻画:
\begin{enumerate}
\renewcommand{\theenumi}{\roman{enumi}}
\renewcommand{\labelenumi}{(\theenumi)}
\item 它对于平展拓扑是局部的.

\item 若 $X'$ 是 $X$ 的 $d$ 个副本的不交和,
从而 $f_{\ast}f^{\ast} F\simeq F^d$, 则迹为求和映射 $F^d\rightarrow F$.
\end{enumerate}

唯一性平凡地由 (i)、(ii) 得到, 因为每个平展覆盖局部常值.
知道唯一性后, 存在性由通常的下降论证得到.

由 (ii) 推出下列公式:
\begin{equation*}
\tr\circ \res:F\longrightarrow F\mbox{ 是乘以 $f$ 的局部次数.}\tag{5.1.4}
\end{equation*}\etiquette[5.1.4]{IX.eq5.1.4}
若次数为常数 $d$, 则态射 $\tr\circ \res: H^q(X,F)\rightarrow H^q(X,F)$
为乘以 $d$.

\begin{enonce*}{推论 5.2}\etiquette[5.2]{IX.cor5.2}
设 $f:X'\rightarrow X$ 为常次数 $d$ 的平展覆盖.
若 $F$ 为 $X$ 上的 $r$-挠层且 $(r,d)=1$,
则 $H^q(X',f^{\ast} F)=0$ 蕴含 $H^q(X,F)=0$.
\end{enonce*}

事实上, 乘以 $d$ 诱导 $F$ 的同构,
从而\pageoriginale 诱导上同调上的同构.
若后一个同构经零分解, 则 $H^q(X,F)=0$.

\setcounter{subsection}{2}
\subsubsection{}\etiquette[5.2.1]{IX.subsubsec5.2.1}
现设 $i:U\rightarrow X$ 为浸入, $f:U'\rightarrow U$ 为次数 $d$ 的平展覆盖,
并假定 $f$ 由一个有限态射 $g:X'\rightarrow X$ 经基变换诱导,
\ie 有笛卡儿图
$$
\xymatrix@C=2cm@R=1.25cm{
U' \ar[r]^{i'}\ar[d]_{f}& X'\ar[d]^{g}\\
U  \ar[r]^{i}& X
}
$$
设 $F$ 为 $U$ 上的层. 有
$$
\xymatrix@C=2.25cm@M=.25cm{
F \ar[r]^{\txt{res}} \ar@/_1.5pc/[rr]_{d}& f_{\ast}f^{\ast}F
\ar[r]^{\txt{tr}}& F,
}
$$
从而, 因 $i_{\ast}f_{\ast}=g_{\ast}i'_{\ast}$,
\hypertarget{sga4.exp.IX.equation.5.3}{}$$
\xymatrix@C=2.25cm{
i_{\ast}F \ar@/_2pc/[rr]_{d} \ar[r]^{\txt{res}}& g_{\ast}i' {}_{\ast}f^{\ast}F \ar[r]^{\txt{tr}}& i_{\ast}F,
}\leqno{(5.3)}
$$
又因 $i_{!}f_{\ast}\simeq g_{\ast}i'_{!}$
(例如由 (VIII \hyperref[VIII.prop5.5]{5.5}) 得到),
$$
\xymatrix@C=2cm{
i_{!}F \ar@/_2pc/[rr]_{d} \ar[r]^{\txt{res}}& g_{\ast}i' {}_{!}f^{\ast}F \ar[r]^{\txt{tr}}& i_{!}F,
}\leqno{(5.4)}
$$\etiquette[5.4]{IX.eq5.4}
从而
$$
\xymatrix@C=2cm{
H^{q}(X,i_{\ast} F) \ar@/_2pc/[rr]_{d'} \ar[r]^{\txt{res}}& H^{q}(X', i' _{\ast}f^{\ast}F) \ar[r]^{\txt{tr}}& H^q(X,i_{\ast}F)
}\leqno{(5.3')}
$$\MakeLinkTarget*{sga4.exp.IX.equation.5.3-prime}\label{5.3'}
且\pageoriginale
$$
\xymatrix@C=2cm{
H^{q}(X,i_{!} F) \ar[r]^{\txt{res}} \ar@/_2.5pc/[rr]_{d}& H^{q}(X', i' _{!}f^{\ast}F)
\ar[r]^{\txt{tr}}& H^q(X,i_{!}F) ;
}\leqno{(5.4')}
$$\MakeLinkTarget*{sga4.exp.IX.equation.5.4-prime}\label{5.4'}
由\eqhyperref[IX.eq5.4]{5.4} 得到下列推论:

\begin{enonce*}{命题 5.5}\etiquette[5.5]{IX.prop5.5}
设 $X$ 为拟紧且拟分离的概形, $\ell \in
\PP$,
$H$ 为取值于 {\rm (Ab)} 的半正合函子,
定义在 $X$ 上 $\ell$-层的范畴上, 且与滤过归纳极限相容.
下列断言等价:
\begin{enumerate}
\renewcommand{\theenumi}{\roman{enumi}}
\renewcommand{\labelenumi}{(\theenumi)}
\item $H=0$.

\item 对每个有限 $f:Y\rightarrow X$,
以及 $Y$ 的每个开子概形 $i:U\rightarrow Y$,
若 $U$ 在 $X$ 上有限表现, 则有 $H(f_{\ast}i_{!}(\ZZ/\ell))=0$.
若 $X$ 诺特, 甚至可限制为整的 $Y$.
\end{enumerate}
\end{enonce*}

证明. 由\hyperref[IX.cor2.7.2]{2.7.2} 及\hyperref[IX.prop2.5]{2.5},
只须对形如 $F=i$, $G$ 的 $F$ 验证 $H(F)=0$,
其中 $i:U\rightarrow X$ 是 $X$ 的一个在 $X$ 上有限表现的子概形的包含,
$G$ 为 $U$ 上局部常值的“等平凡”$\ell$-层.
设 $U''\rightarrow U$ 为群为 $g$ 的主覆盖,
使得 $g$ 在 $U''$ 上变为常值的.
令 $A$ 为 $G$ 在 $U''$ 上的值所成的 $\ell$-群.
群 $g$ 在 $A$ 上作用, 这个作用决定 $G$ 的结构.
取 $g$ 的一个 $\ell$-Sylow 子群 $\sheaf{h}$,
令 $U'=U''/\sheaf{h}$ 为 $U$ 的对应于 $\sheaf{h}$ 的覆盖.
$U'$ 在 $U$ 上的次数 $d$ 与 $\ell$ 互素.
选取 $f:U'\rightarrow U$ 到一个有限态射 $g:X'\rightarrow X$ 的延拓,
这样的延拓总存在 (\EGA IV 8.12.6),
令 $i':U'\rightarrow X'$ 为包含.
应用\eqhyperref[IX.eq5.4]{5.4} 及 $d$ 与 $\ell$ 互素这一事实,
归结为证明 $H(g_{\ast}i'_{!}f^{\ast} G)=0$.

现在\pageoriginale 众所周知,
在一个 $\ell$-群 $\sheaf{h}$ 的作用下为单的阿贝尔 $\ell$-挠群,
只能是带有平凡作用的 $\ZZ/\ell$;
所以 $A$ 作为 $\sheaf{h}$-模有一个滤过,
其相继商均同构于带有平凡作用的 $\ZZ/\ell$.
这一滤过给出 $f^{\ast} G$ 的相应滤过,
因此要证明 $H(g_{\ast}i'_{!} f^{\ast} G)=0$,
只须证明 $H(g_{\ast}i'_{!}\ZZ/r)=0$, 由此即得结果.

\begin{enonce*}{命题 5.6}\etiquette[5.6]{IX.prop5.6}
设 $X$ 为诺特概形, $\ell \in \PP$,
$\varphi$ 为定义在 $X$ 上有限的整概形的集合上、取值于 $\NN$ 的函数.
假定只要存在 $X$-态射 $Y_1\rightarrow Y_2$,
就有 $\varphi(Y_1)\leq \varphi(Y_2)$,
并且若 $Y_1$ 是 $Y_2$ 的一个不同于 $Y_2$ 的闭子概形, 则不等式严格.
对 $X$ 上每个有限 $Y$, 置 $(Y)=\Sup (Y_i)$,
其中各 $Y_i$ 为 $Y$ 的不可约分支, 带诱导的既约结构.
则下列断言等价:
\begin{enumerate}
\renewcommand{\theenumi}{\roman{enumi}}
\renewcommand{\labelenumi}{(\theenumi)}
\item 对 $X$ 上每个有限 $Y$ 及 $Y$ 上每个 $\ell$-层 $F$,
当 $q>\varphi(\Supp F)$ 时有 $H^q(Y,F)=0$.

\item 对 $X$ 上每个有限的整概形 $Y$,
当 $q>\varphi(Y)$ 时有 $H^q(Y,\ZZ/\ell)=0$.
\end{enumerate}
\end{enonce*}

证明. 假定 (ii) 成立.
可以对每个 $H^q$ 应用上面的命题,
考虑到上同调与归纳极限以及有限态射下的正像可交换;
因此要验证 (i), 归结到 $Y$ 为 $X$ 上有限的整概形,
且 $F$ 形如 $i_{!}(\ZZ/\ell)_U$ 的情形,
其中 $i:U\rightarrow X$ 是非空开子概形.

有正合列
$$
0\longrightarrow i_{!}(\ZZ/\ell)_U\longrightarrow
(\ZZ/\ell)_Y\longrightarrow
(\ZZ/\ell)_Z\longrightarrow 0
$$
其中 $Z=Y-U$, 所以 $Z\neq Y$.
对 $\varphi$ 的值作归纳, 可\pageoriginale 假定当 $q>\varphi (Y)-1$ 时,
$H^q(\ZZ, \ZZ/\ell)=0$,
从而由上同调正合列, 当 $q>\varphi(Y)$ 时有 $H^q(Y,i_{!}(\ZZ/\ell))=0$.

\begin{enonce*}{推论 5.7}\etiquette[5.7]{IX.cor5.7}
设 $X$ 为可分闭域 $k$ 上的代数曲线,
其特征不同于 $\ell$, 其中 $\ell\in \PP$.
则 $X$ 以 $\ell$-挠层 $F$ 为系数的上同调在次数 $>2$ 时为零.
若 $X$ 仿射, 则上同调在次数 $>1$ 时为零.
\end{enonce*}

对 $X$ 上有限的整概形 $Y$, 置 $\varphi(Y)=2\dim Y$
(\resp 若 $X$, 因而 $Y$ 仿射, 则置 $\varphi(Y)=\dim Y$).
由\hyperref[IX.prop5.6]{7.6},
只须证明当 $q>\varphi(Y)$ 时有 $H^q(Y,\ZZ/\ell)=0$;
若 $\dim Y=0$, 这平凡成立.
所以只须对整曲线 $Y$, 证明当 $q>2$
(\resp 仿射情形下当 $q>1$) 时有 $H^q(Y,\ZZ/\ell)=0$.
因为 $(\ZZ/\ell)_Y\simeq (𝛍_{\ell})_Y$,
对 $q>2$ 这由\hyperref[IX.thm4.6]{4.6} 得到.

设 $X$ 仿射.
因为根式扩张不影响上同调 (VIII \hyperref[VIII.thm1.1]{1.1}),
归结到 $k$ 代数闭的情形.
设 $f:X'\rightarrow X$ 为 $X$ 的正规化,
它是有限态射, 且在有限点集之外为同构, 将该有限点集记为 $S$.
在 $X$ 上有层的正合列
$$
0\longrightarrow (\ZZ/n)_X\longrightarrow
f_{\ast}(\ZZ/n)_{X'}\longrightarrow \epsilon
\longrightarrow 0{\quoi},
$$
其中 $\epsilon$ 集中在 $S$ 上,
所以当 $q>0$ 时有 $H^q(X,\epsilon)=0$.
由 (VIII \hyperref[VIII.cor5.6]{5.6}), 归结为证明当 $q>1$ 时,
$$
H^q(X,f_{\ast}(\ZZ/n)_{X'})\simeq
H^q(X',\ZZ/n)=0
$$
, 即归结到 $X$ 正规的情形.
显然还可设 $X$ 连通.
取开浸入 $i:X\rightarrow \bar{X}$,
其中 $\bar{X}$ 完备且非奇异.
因为 $X$ 仿射, $\bar{X}$ 中至少缺少一个点,
由此立即得知态射
$$
A(k)\longrightarrow \Pic X
$$
为\pageoriginale 满射, 其中 $A$ 为 $\bar{X}$ 的 Jacobian
(\cf \eqhyperref[IX.eq4.8]{4.8}).
因为 $A(k)$ 可被 $\ell\neq \car k$ 除,
$\Pic X$ 也如此,
从而由\hyperref[IX.thm4.6]{4.6}, 有 $H^2(X,𝛍_{\ell})=0$, \cqfd

\begin{enonce*}[remark]{注 5.8}\etiquette[5.8]{IX.rem5.8}
为便于以后引用, 记录下\hyperref[IX.prop5.5]{5.5} 的论证实际上确立的结果:
设 $X$ 为拟紧且拟分离的概形, $\ell$ 为素数,
$C$ 为可构造阿贝尔 $\ell$-挠层范畴的一个严格全子范畴,
对直因子及扩张稳定.
假定对每个有限态射 $f:Y\rightarrow X$
及 $Y$ 的每个在 $X$ 上有限表现的开子概形 $i:U\rightarrow Y$,
都有 $f_{\ast}(i_{!}(𝐙/\ell𝐙)_U)\in C$.
则 $C$ 包含所有可构造的
%% 𝒵 here is ℤ [check the macros]
$\ell$-挠层.
\end{enonce*}

\begin{thebibliography}{99}
\bibitem{sga4.exp.IX.bib.1} Gabriel, P. Des Catégories Abéliennes,
学位论文, Paris (1962).

\bibitem{sga4.exp.IX.bib.2} Ogg, A. Cohomology of abelian varieties over
functions fields, Annals of Math. Vol. 76, No. 2 (1962).

\bibitem{sga4.exp.IX.bib.3} Deuring, M. Die Typen der Multiplikatorenringe elliptischer
Funktionenkörper, Hamb. Abh. Bol. 13 (1940) p. 197.

\bibitem{sga4.exp.IX.bib.4} Grothendieck, A. Fondements de la Géométrie
Algébrique, Extraits du Séminaire Bourbaki 1957-1962), Secrétariat
mathématique, 11, rue Pierre Curie, Paris 5ème. (\emph{引作} FGA).

\bibitem{sga4.exp.IX.bib.5} Serre, J.P. Cohomologie Galoisienne, Lecture Notes
nᵒ 5, (Springer) (\emph{引作} CG).

\bibitem{sga4.exp.IX.bib.6} Douady, A. Cohomologie des groupes compacts
totalement discontinus, Séminaire Bourbaki, exp. 189.
\end{thebibliography}

\ifx\danslelivre\undefined
\end{document}
\fi

% END INCLUDED SGA4_English_current_buildable_source_20260906/source/09/drafts/SGA4_IX_LINES_1607_1883_English.texfrag


